\documentclass[11pt,reqno]{amsart}
\usepackage[T1]{fontenc}
\usepackage{lmodern}
\usepackage[a4paper,margin=28mm,headheight=14pt]{geometry}
\usepackage{amsmath,amssymb,mathtools,booktabs,longtable,array,microtype}
\usepackage[hidelinks,unicode]{hyperref}
\usepackage{url}
\emergencystretch=2em
\allowdisplaybreaks[2]
\numberwithin{equation}{section}
\newtheorem{theorem}{Theorem}[section]
\newtheorem{lemma}[theorem]{Lemma}
\newtheorem{proposition}[theorem]{Proposition}
\newtheorem{corollary}[theorem]{Corollary}
\newtheorem{certificate}[theorem]{Finite certificate}
\newtheorem{reported}[theorem]{Source-qualified statement}
\theoremstyle{definition}
\newtheorem{definition}[theorem]{Definition}
\newtheorem{assumption}[theorem]{Assumption}
\newtheorem{problem}[theorem]{Open problem}
\theoremstyle{remark}
\newtheorem{remark}[theorem]{Remark}
\newcommand{\N}{\mathbb N}
\newcommand{\Q}{\mathbb Q}
\newcommand{\Z}{\mathbb Z}
\newcommand{\R}{\mathbb R}
\newcommand{\C}{\mathbb C}
\newcommand{\E}{\mathbb E}
\newcommand{\eps}{\varepsilon}
\newcommand{\ind}{\mathbf 1}
\newcommand{\ii}{\mathrm i}
\newcommand{\dd}{\,\mathrm d}
\newcommand{\abs}[1]{\lvert#1\rvert}
\newcommand{\norm}[1]{\lVert#1\rVert}
\DeclareMathOperator{\rad}{rad}
\DeclareMathOperator{\Tr}{Tr}
\DeclareMathOperator{\Res}{Res}
\DeclareMathOperator{\Var}{Var}
\DeclareMathOperator{\cond}{cond}
\DeclareMathOperator{\arcosh}{arcosh}
\DeclareMathOperator{\arsinh}{arsinh}
\renewcommand{\Re}{\operatorname{Re}}
\renewcommand{\Im}{\operatorname{Im}}
\newcommand{\git}[2]{\href{https://github.com/GettysburgResearch/riemann/blob/#1/#2}{\texttt{#1}}}
\title[Arithmetic transfers and finite theta approximants]{Arithmetic transfers, finite theta approximants,\texorpdfstring{\\}{ }and obstruction estimates for zeta functions}
\author{Gideon Freund}
\address{Gettysburg Research}
\date{19 October 2026}
\subjclass[2020]{Primary 11M26; Secondary 11M06, 11M41, 11N37, 30D15, 82B20}
\keywords{Riemann zeta function, quadratic characters, M\"obius function, mollifiers, theta moments, canonical products, determinant regularization}
\thanks{This draft consolidates research and computational material developed with substantial AI assistance. Its provenance, verification scope, and remaining author-review obligations are recorded in Section~\ref{sec:verification} and the accompanying source ledger.}
\hypersetup{
 pdftitle={Arithmetic transfers, finite theta approximants, and obstruction estimates for zeta functions},
 pdfauthor={Gideon Freund},
 pdfsubject={Source-grounded mathematical preprint draft; no RH proof or new ordinary zero-free region claimed},
 pdfkeywords={zeta function, quadratic characters, Mobius, theta moments, canonical products}
}
\begin{document}
\begin{abstract}
We present a source-grounded account of recent arithmetic and approximation results in the Gettysburg Research Riemann repository. The common objective is to preserve the contribution of a fixed zeta zero while changing the arithmetic representation in which it must be estimated. We give a conductor-uniform comparison between primitive quadratic-field reciprocal coefficients and ordinary M\"obius coefficients, with an explicit finite-family error. For actual quadratic characters, a mean-square approximation above the critical line yields quantitative full-residual transfer and an exact optimization that removes auxiliary local variance without removing the ordinary zeta factor. We also present a two-mollifier transfer theorem for capped reciprocal moments, complete native denominator and frequency-sector bounds, and positive cycle determinants with explicitly controlled regularization. On the approximation side, the frozen computational sources provide five directed certificates of real-zero entire functions matching the normalized theta expansion through orders 22, 24, and 26. A positive weighted canonical-moment criterion removes several unnecessary restrictions from the remaining all-order construction. We explain a compactness obstruction for strip-confined probability laws, give a corrected strip-transfer theorem and a finite-graph extension admitting controlled negative couplings, and preserve an exact failed comparison for the proposed boundary $0.87$. The manuscript distinguishes proved component arguments, source-qualified analytic deductions, finite certificates, and unproved sufficient estimates. It claims neither the Riemann hypothesis nor a new global zero-free region for the ordinary zeta function.
\end{abstract}
\maketitle
\tableofcontents
\section{Introduction and scope}\label{sec:introduction}

An averaged estimate, a finite real-zero construction, and a bound for a fixed zeta function are different mathematical statements. Much of the recent work considered here concerns the operations needed to pass between them without discarding the original arithmetic source. The point of this manuscript is to state the resulting contributions in conventional mathematical form, with their dependencies and quantifiers intact.

Our source window is the research developed on 17--19 October 2026, together with the earlier arguments on which those packets explicitly depend. The principal sources are the frozen contributions \cite{R977,R989,R993,R994,R996,R998,R1000,R1001,R1002,R1003,R1004}. These are research-branch publications, not a single integrated proof. In particular, the latest source \cite{R1003} supersedes the earlier order-20 frontier for \emph{finite Taylor matching}, but does not supply an unbounded construction. Recoveries of unchanged files are not counted as additional discoveries.

\subsection{The contributions in this account}
The arithmetic part begins with a primitive rational sequence associated with an imaginary quadratic field. Its local factors distinguish split, inert, and ramified primes. An absolutely summable invertible correction compares this sequence with the complete reciprocal Dedekind source. Averaging over actual fundamental discriminants then gives an explicit approximation to ordinary M\"obius, uniformly for bounded common weights. The comparison preserves the common M\"obius term rather than using character orthogonality to remove it.

A second family argument concerns analytically continued quadratic $L$-values on fixed compact subsets of $1/2<\Re s<1$. Starting from the classical quadratic large sieve, one retains square parts, conductor-dependent weights, and the complete dual expression in a flat-cutoff functional equation. This gives a quantitative transfer of full mollified residuals to a random Euler model. A finite-state minimization then identifies the optimal local correction at a prescribed mean. The auxiliary variance tends to zero, but the rational residual $\abs{\zeta(s)M_Z(s)-1}^2$ remains.

The two-mollifier argument has a different purpose: one polynomial controls large reciprocal values while another computes the leading term. Its abstract transfer theorem is independent of the deep character-family inputs. Applications to every fixed reciprocal moment order retain an essential cap. The cap, its size, and the loss of uniformity as the moment order grows are part of the conclusion, not removable technicalities.

For the native rational M\"obius source, the complete denominator decomposition pays sectors that earlier product cutoffs retained. The small-core estimate includes all coherent cross terms and arbitrary bounded denominator-dependent masks. A separate Fourier estimate removes the entire high-frequency tail at a specified power cutoff. These estimates leave a concrete low-frequency or large-core arithmetic remainder, rather than an unspecified error term.

The approximation part has both finite and all-order components. The sources provide directed existence certificates for functions
\[
 P(z)=e^{-az^2}\prod_{j=1}^N(1-b_jz^2),\qquad a,b_j>0,
\]
matching the actual standardized theta expansion through order 26. They also establish a conditional canonical-product theorem with positive \emph{real weights}, not only integer root multiplicities. A separate strip programme admits finite Ising models with controlled fields and negative couplings. None of these all-order criteria is supplied with its unbounded feasibility premise.

Finally, positive primitive-cycle determinants realize the exact integer
\[
 A_N=\prod_{p^j\le N}p^{N-p^j}.
\]
Their regularization error is explicitly bounded, including a larger-mass regime with improved conditioning. The missing estimate is the centered logarithmic height, not positive definiteness. An exact exponent audit records why the attempted $0.87$ comparison did not follow from the available sufficient estimates.

\subsection{Claim discipline and the role of the sources}
The self-contained arguments below are written as lemmas and theorems. Their inclusion records the mathematical claims and proofs being offered for review; it is not an independent referee verdict. Results involving deep analytic inputs state those inputs explicitly. The finite certificates are identified by their exact source data and acceptance rules; the compilation of this manuscript is not described as a fresh execution of every source checker. Source-qualified zero-free proposals are separated in Section~\ref{sec:boundaries}, and are not premises of the elementary quadratic comparison, canonical-factor, or determinant arguments.

The manuscript does not silently repair a source, infer a missing analytic estimate from a finite test, or identify a new representation with a solution of its remaining bound. In particular, neither an $L^2$ average over many characters nor a finite sequence of matching orders rules out a prescribed exceptional zero. We make no claim of exhaustive novelty against the literature. Classical ingredients are credited where they enter; the proposed repository contributions consist of the displayed constructions, quantitative comparisons, scoped applications, and explicit obstruction calculations.

\subsection{Notation and elementary zero retention}
We write $\mu$ and $\Lambda$ for the M\"obius and von Mangoldt functions, and
\[
 M(x)=\sum_{n\le x}\mu(n),\qquad
 m(x)=\sum_{n\le x}\frac{\mu(n)}n,\qquad
 F_X=\sum_{k\le X}\abs{m(k)}^2.
\]
The harmonic number is $h_N=\sum_{j\le N}j^{-1}$. For a number field, $N\mathfrak a$ denotes the absolute ideal norm. Superscript $S$ denotes removal of the stated finite Euler factors. Every occurrence of $O_\eps$ or $\ll_\eps$ allows dependence on the explicitly fixed data. Moment orders are fixed unless uniformity is expressly stated.

For a quadratic discriminant $\Delta$, the factorization
\begin{equation}\label{eq:quadratic-factorization}
 \zeta_{\Q(\sqrt\Delta)}(s)=\zeta(s)L(s,\chi_\Delta)
\end{equation}
is used in its usual meromorphic sense; here and below $\Q$ denotes $\mathbb Q$. The nonprincipal Dirichlet factor is entire. Thus a nontrivial ordinary zeta zero is a zero of every such full Dedekind product. This observation is lost if one instead studies only the quotient by $\zeta$.

For Fourier arguments we use
\[
 \Xi(z)=\xi(1/2+\ii z),\qquad
 \xi(s)=\tfrac12s(s-1)\pi^{-s/2}\Gamma(s/2)\zeta(s).
\]
A zero $\rho=\beta+\ii\gamma$ corresponds to $z=\gamma-\ii(\beta-1/2)$. Hence confinement of all zeros of $\Xi$ to $\abs{\Im z}\le b$ implies $\Re\rho\le 1/2+b$. A standardized theta variable changes this band by its scale factor; Section~\ref{sec:theta} fixes that normalization explicitly.

Two elementary warnings recur. If $L(\rho)=0$ and $M$ is finite at $\rho$, then $L(\rho)M(\rho)-1=-1$. Its squared residual contributes one, not infinity. Conversely, if a fixed-source Mellin transform has a nonremovable pole at every target zero, proving holomorphy of that transform in a region really does exclude those zeros. The issue is to prove the required bound on the same source.

\section{Quadratic-field coefficients and a uniform signed average}\label{sec:quadratic}

A quadratic field has a useful feature: its zeta function contains ordinary $\zeta$ as a factor. Varying the field therefore gives many different arithmetic expressions with the same target factor. The question is whether averaging these expressions can simplify the auxiliary arithmetic without losing ordinary $\zeta$.

The first result replaces ideal coefficients by an explicit sequence on rational integers. The replacement preserves any cancellation estimate with exponent greater than $1/2$. The second result shows quantitatively that averaging this sequence over quadratic fields recovers the ordinary M\"obius sequence. Neither step proves that the resulting sum is small; it tells us exactly which sum must be estimated.

The fixed-field construction is from \cite{R996}, and the uniform extension and field average are from \cite{R1002}. The family here differs from the dyadic family in Section~\ref{sec:residuals}; we specify each one rather than interchange them.

\subsection{Replacing ideal coefficients by rational coefficients}
Put $\mu_6(n)=\mu(n)\ind_{(n,6)=1}$. For a negative fundamental discriminant $\Delta$, let $\chi_\Delta$ be the associated primitive quadratic character. Define
\begin{equation}\label{eq:primitive-sequences}
 a_\Delta=\mu_6*(\mu_6\chi_\Delta),\qquad
 b_\Delta(n)=\mu_6(n)\prod_{p\mid n}(1+\chi_\Delta(p)),
\end{equation}
where $*$ denotes Dirichlet convolution. Thus $b_\Delta$ is supported on squarefree integers coprime to $6$. In particular a ramified prime $p\ge5$ has coefficient $-1$, not zero. Write $\zeta_6$ and $L_6$ for deletion of the Euler factors at $2,3$.

\begin{theorem}[Uniform primitive-core comparison]\label{thm:primitive}
On $\Re s>1$, let
\[
 A_\Delta(s)=\sum_n a_\Delta(n)n^{-s}
 =\frac1{\zeta_6(s)L_6(s,\chi_\Delta)},\qquad
 B_\Delta(s)=\sum_n b_\Delta(n)n^{-s}.
\]
There are Dirichlet series $V_\Delta,W_\Delta$ such that
\begin{equation}\label{eq:primitive-inverse}
 B_\Delta=A_\Delta V_\Delta,\qquad
 A_\Delta=B_\Delta W_\Delta,\qquad V_\Delta W_\Delta=1.
\end{equation}
For every fixed $\sigma>1/2$, their coefficients satisfy
\[
 \sum_n\frac{\abs{v_\Delta(n)}}{n^\sigma}\le C_V(\sigma),\qquad
 \sum_n\frac{\abs{w_\Delta(n)}}{n^\sigma}\le C_W(\sigma),
\]
with constants independent of $\Delta$. Both functions are holomorphic and nonvanishing on $\Re s>1/2$.
\end{theorem}
\begin{proof}
At $p\ge5$ put $x=p^{-s}$. At a split prime the local factors are
\begin{align*}
 A_p(x)&=(1-x)^2,& B_p(x)&=1-2x,\\
 V_p(x)&=\frac{1-2x}{(1-x)^2}
       =1-\sum_{j\ge2}(j-1)x^j,&
 W_p(x)&=\frac{(1-x)^2}{1-2x}
       =1+\sum_{j\ge2}2^{j-2}x^j.
\end{align*}
At an inert prime they are
\[
 A_p=1-x^2,\quad B_p=1,\quad
 V_p=(1-x^2)^{-1},\quad W_p=1-x^2.
\]
At a ramified prime $A_p=B_p=1-x$, so $V_p=W_p=1$. All factors at $2,3$ are one. These are identities of complete power series; the higher prime powers are not discarded.

Since $2p^{-\sigma}<1$ for $p\ge5$ and $\sigma>1/2$, the products
\begin{align*}
 C_V(\sigma)&=\prod_{p\ge5}
  \left(1+\frac{p^{-2\sigma}}{(1-p^{-\sigma})^2}\right),\\
 C_W(\sigma)&=\prod_{p\ge5}
  \left(1+\frac{p^{-2\sigma}}{1-2p^{-\sigma}}\right)
\end{align*}
converge. They dominate the corresponding local absolute coefficient sums in every splitting case. This proves the uniform bounds and normal convergence on smaller half-planes. The local inverse identities pass to the products, proving nonvanishing and \eqref{eq:primitive-inverse}.
\end{proof}

\begin{corollary}[No loss in a fixed power exponent]\label{cor:primitive-power}
Fix $\beta>1/2$. A bound
\[
 \sum_{n\le X}a_\Delta(n)\ll_{\Delta,\eps}X^{\beta+\eps}
\]
for every $\eps>0$ is equivalent to the same bound with $b_\Delta$ in place of $a_\Delta$. The same equivalence holds for estimates uniform over additive phases, or over all real polynomial phases of any one fixed degree.
\end{corollary}
\begin{proof}
For example, $b_\Delta=a_\Delta*v_\Delta$ gives
\[
 \sum_{n\le X}b_\Delta(n)
 =\sum_{d\le X}v_\Delta(d)\sum_{m\le X/d}a_\Delta(m).
\]
Apply the hypothesized bound with a smaller positive loss and sum $\abs{v_\Delta(d)}d^{-\beta-\eps'}$. The inverse uses $W_\Delta$. For a polynomial phase $e^{2\pi\ii P(n)}$, the inner phase becomes $P(dm)$, which is still a polynomial of the same degree. The equivalence transfers an estimate; it does not establish one.
\end{proof}

At the all-$\eps$ square-root exponent, this reproduces the fixed-field reciprocal criterion described in \cite{R996}. For $K=\Q(\sqrt{-3})$, the unramified primitive sequence is simply $\mu(n)2^{\omega(n)}$ on squarefree integers all of whose prime factors are $1\bmod3$. Its fixed-field square-root estimate is stronger than ordinary RH alone, because it also controls the factor $L(s,\chi_{-3})$.

\subsection{Averaging over quadratic fields}
For $Q\ge2^{16}$ define
\[
 \mathcal D(Q)=\{d\le Q:d>0,\ d\text{ squarefree},\ d\equiv3\pmod4\},
 \qquad N_Q=\abs{\mathcal D(Q)}.
\]
Every $-d$ is a fundamental discriminant. We average uniformly over this complete set.

\begin{lemma}\label{lem:quadratic-mean}
One has $N_Q\ge Q/8$. For every odd squarefree $r>1$,
\begin{align}
 \left|\frac1{N_Q}\sum_{d\in\mathcal D(Q)}\chi_{-d}(r)\right|
 &\le24\sqrt{r/Q},\label{eq:mean-trivial}\\
 \left|\frac1{N_Q}\sum_{d\in\mathcal D(Q)}\chi_{-d}(r)\right|
 &\le\frac{24r^{1/4}\sqrt{1+\log r}}{\sqrt Q}.
 \label{eq:mean-pv}
\end{align}
\end{lemma}
\begin{proof}
There are at least $Q/4-1$ integers $3\bmod4$ up to $Q$. For each odd prime $p\le\sqrt Q$, at most $Q/(4p^2)+1$ are divisible by $p^2$. Using $\sum_{p\ge3}p^{-2}\le\sum_{n\ge3}n^{-2}\le4/9$ gives
\[
 N_Q\ge\frac{5Q}{36}-\sqrt Q-1\ge\frac Q8.
\]

Expand $\mu^2(d)=\sum_{a^2\mid d}\mu(a)$. Terms with $(a,r)>1$ vanish and are retained as zero terms. For the remaining odd $a$, the inner character sum is over a progression $4j+3$ of length at most $Q/a^2$. The Jacobi character modulo the odd squarefree integer $r$ is primitive and nonprincipal. Its complete mean is zero, so the inner sum is at most $\min(r,Q/a^2)$ in modulus. More generally, any bound $T$ for its incomplete sums gives
\[
 \sum_{a\le\sqrt Q}\min(T,Q/a^2)\le3\sqrt{QT}.
\]
This follows by splitting at $\sqrt{Q/T}$; when that point is below one, the bound follows directly from $\sum a^{-2}<2$ and $T>Q$.

Taking $T=r$ and dividing by $N_Q$ proves \eqref{eq:mean-trivial}. For \eqref{eq:mean-pv}, primitive Gauss-sum Fourier inversion and
$\abs{1-e^{2\pi\ii a/r}}\ge4\norm{a/r}_{\R/\Z}$ give the elementary P\'olya--Vinogradov bound
\[
 T\le\sqrt r\,h_{\lfloor r/2\rfloor}
 \le\sqrt r(1+\log r).
\]
The fixed invertible change of residue induced by $4j+3$ has no effect on that bound. Substituting this $T$ completes the proof. No ramified value of a character has been replaced by a unit-modulus value.
\end{proof}

\begin{theorem}[Uniform weighted finite-family comparison]\label{thm:field-average}
For $X\ge1$, $Q\ge2^{16}$, and every common complex array $\abs{u_n}\le1$,
\begin{align}
 &\left|\frac1{N_Q}\sum_{d\in\mathcal D(Q)}
       \sum_{n\le X}b_{-d}(n)u_n
       -\sum_{n\le X}\mu_6(n)u_n\right|
 \le\frac{48X^{3/2}}{\sqrt Q},\label{eq:field-average-weak}\\
 &\left|\frac1{N_Q}\sum_{d\in\mathcal D(Q)}
       \sum_{n\le X}b_{-d}(n)u_n
       -\sum_{n\le X}\mu_6(n)u_n\right|
 \le\frac{96X^{5/4}\sqrt{1+\log X}}{\sqrt Q}.
 \label{eq:field-average}
\end{align}
\end{theorem}
\begin{proof}
For squarefree $n$,
\[
 \prod_{p\mid n}(1+\chi_{-d}(p))=\sum_{r\mid n}\chi_{-d}(r).
\]
The term $r=1$ is exactly the common M\"obius term in the theorem. It is canceled algebraically before taking absolute values. Every other divisor is squarefree and odd, so \eqref{eq:mean-trivial} bounds the remaining expression by
\[
 \frac{24}{\sqrt Q}\sum_{n\le X}\sum_{\substack{r\mid n\\r>1}}\sqrt r
 \le\frac{24X}{\sqrt Q}\sum_{r\le X}r^{-1/2}
 \le\frac{48X^{3/2}}{\sqrt Q}.
\]
Using \eqref{eq:mean-pv} instead gives
\[
 \frac{24X\sqrt{1+\log X}}{\sqrt Q}\sum_{r\le X}r^{-3/4}
 \le\frac{96X^{5/4}\sqrt{1+\log X}}{\sqrt Q}.
\]
Squarefreeness and coprimality were relaxed only in these positive upper sums, not in the original observable.
\end{proof}

Theorem~\ref{thm:field-average} is a quantitative arithmetic comparison, not a cancellation estimate for its first term. If $\theta>1/2$ and $Q=X^q$ with $q\ge5/2-2\theta$, its error fits $X^{\theta+\eps}$. Thus a signed average bound at that power would imply the same bound for $\sum_{n\le X}\mu_6(n)$, and hence ordinary zero exclusion to the right of $\theta$ by Abel summation. The schedules $Q=X^{5/6}$ and $Q=X^{3/2}$ correspond, respectively, to the still-open $5/6$ and square-root targets. They are not achieved zero-free boundaries.

The order of this deduction matters. The finite-family comparison is first applied at each $X$ to recover a \emph{fixed} coefficient sequence. One does not take a Mellin transform while treating a family whose size depends on $X$ as though it were a fixed Dirichlet series. Similarly, arbitrary signed combinations of individual reciprocal residues need not preserve those residues. The theorem supplies a specific coefficient comparison, not a blanket noncancellation assertion.

Finally, deletion of an exceptional set must be paid. Uniformly in $d$,
\[
 \left|\sum_{n\le X}b_{-d}(n)u_n\right|
 \le\sum_{n\le X}2^{\omega(n)}\le X(1+\log X).
\]
Removing $E$ rows while retaining normalization $N_Q$ therefore costs at most $(E/N_Q)X(1+\log X)$. At $Q=X^{5/6}$, an exceptional proportion $Q^{-\delta}$ requires $\delta\ge1/5$ merely to fit an $X^{5/6+\eps}$ allowance. An assertion about almost every auxiliary factor is not a bound for the shared ordinary source.

\section{Actual quadratic families and optimal auxiliary-variance removal}\label{sec:residuals}

We now use the dyadic family from \cite{R1001,R1003},
\[
 \mathcal U(Q)=\{u:Q<u\le2Q,\ u\text{ odd and squarefree}\},
 \qquad \chi_u=\chi_{-8u}.
\]
These are primitive odd characters of conductor $q=8u$ and root number one. Write $\E_Q$ for the uniform family average. The cardinality is $(4/\pi^2)Q+O(\sqrt Q)$. We retain all ramified values.

\subsection{The character mean and the random square kernel}
The elementary squarefree-character calculation in \cite[shared proof]{R1001}, restated in \cite[arithmetic proof, (1.1)]{R1003}, is
\begin{equation}\label{eq:dyadic-character-mean}
 \E_Q\chi_u(n)=\ind_{n=\square}g(n)
   +O\bigl(Q^{-1/2}n^{1/4}\log(2n)\bigr),\qquad
 g(n)=\prod_{p\mid n}\frac p{p+1},
\end{equation}
for odd $n$. For even $n$ both the actual and limiting characters vanish. The calculation expands the squarefree indicator and applies primitive character completion below a square-divisor cutoff; the square case retains the coprimality probability $g(n)$. We use precisely this stated finite-family estimate, not a zero-density assertion.

Let the independent variables at odd primes have distribution
\begin{equation}\label{eq:local-law}
 \Pr(X_p=0)=\frac1{p+1},\qquad
 \Pr(X_p=1)=\Pr(X_p=-1)=\frac p{2(p+1)},\qquad X_2=0,
\end{equation}
and extend them completely multiplicatively. Then $\E X(n)=\ind_{n=\square}g(n)$ for odd $n$. Prescribing finitely many local states is permitted, provided their limiting probability is positive. The indicator of one such prescription is a finite product of the polynomials $1-X_p^2$ and $(X_p^2\pm X_p)/2$. Expanding it in \eqref{eq:dyadic-character-mean} proves the same mean estimate under conditioning, with a constant depending on the fixed set of primes. Neither a growing set of prescribed states nor uniformity as its probability tends to zero is asserted.

A useful exact consequence concerns the inverse Dedekind coefficients. Define
\[
 a_u=\mu*(\mu\chi_u),\qquad A_u(x)=\sum_{n\le x}a_u(n).
\]
For fixed $x,y$, expanding the finite sums in \eqref{eq:dyadic-character-mean} gives
\begin{align}
 \lim_{Q\to\infty}\E_Q A_u(x)&=M(x),\label{eq:mean-Mertens}\\
 \lim_{Q\to\infty}\E_Q[A_u(x)A_u(y)]
 &=\sum_{\substack{m\le\min(x,y)\\m\text{ odd squarefree}}}
       g(m)M(x/m)M(y/m).\label{eq:cov-Mertens}
\end{align}
Indeed $A_u(x)=\sum_{m\le x}\mu(m)\chi_u(m)M(x/m)$, and two squarefree indices have square product exactly when they are equal. The term $m=1$ is the original ordinary covariance. The additional diagonal weights are ramification probabilities, not independent copies of the original source. These fixed-cutoff limits alone do not authorize a growing-cutoff interchange.

\subsection{Mean-square approximation of the complete \texorpdfstring{$L$}{L}-value}
Fix a compact set $\mathcal K\subset\{s:1/2<\Re s<1\}$. All heights below are bounded by this compact set. Write
\[
 \sigma=\Re s,\qquad h=\sigma-\tfrac12,\qquad
 \lambda=\tfrac54-\sigma,\qquad \sigma_0=\min_{s\in\mathcal K}\Re s.
\]
The external analytic input is Heath-Brown's quadratic large sieve \cite{HB95}; see also the statement in \cite[Theorem 2.4]{SY10}:
\begin{equation}\label{eq:quadratic-LS}
 \sum_{\substack{u\le U\\u\text{ odd squarefree}}}
 \left|\sum_{\substack{a\le V\\a\text{ odd squarefree}}}
        c_a\left(\frac au\right)\right|^2
 \ll_\eps (UV)^\eps(U+V)\sum_a\abs{c_a}^2.
\end{equation}
Quadratic reciprocity and fixed residue-class decompositions give the corresponding inequality for $\chi_u(a)$.

\begin{lemma}[All square parts retained]\label{lem:full-column-LS}
Let $V\ge1$, $\abs{c_n}\le n^{-\beta}$ on $V<n\le2V$, with $\beta$ in a fixed bounded interval. Suppose the row-dependent weights $W_u(n)$ have uniformly bounded supremum and total variation on this interval. Then
\begin{equation}\label{eq:full-column-LS}
 \E_Q\left|\sum_{V<n\le2V}c_n\chi_u(n)W_u(n)\right|^2
 \ll_\eps (QV)^\eps(1+V/Q)V^{1-2\beta}.
\end{equation}
A partial dyadic endpoint and any fixed positive-density local condition are allowed.
\end{lemma}
\begin{proof}
Write uniquely $n=ab^2$, with $a$ odd squarefree and $b$ odd; there is no condition $(a,b)=1$. Exactly,
\[
 \chi_u(ab^2)=\chi_u(a)\ind_{(u,b)=1}.
\]
For each fixed $b$, keep the row restriction before enlarging the nonnegative row sum in \eqref{eq:quadratic-LS}.

The maximal initial-segment form of that inequality costs at most two logarithmic factors: decompose each initial segment into binary intervals, apply Cauchy--Schwarz on each row, and then sum the large-sieve estimates for all binary intervals. A coefficient is counted at most logarithmically many times. Abel summation now permits the stated row-dependent weights. For fixed $b$, the coefficient square mass is $O(V^{1-2\beta}/b^2)$, including a last interval with $V/b^2$ between $1/2$ and $1$. Its normalized row norm is at most
\[
 (QV)^\eps(1+V/Q)^{1/2}V^{1/2-\beta}/b.
\]
Minkowski and $\sum_{b\le\sqrt{2V}}b^{-1}\ll\log(2V)$ prove the assertion after shrinking the preliminary losses. The harmonic square-part cost has not been replaced by a convergent sum.
\end{proof}

\begin{theorem}[Full-$L$ tail estimate]\label{thm:L-tail}
For $1\le K\le Q$, $s\in\mathcal K$, and every $\eps>0$,
\begin{equation}\label{eq:L-tail}
 \E_Q\left|L(s,\chi_u)-\sum_{n\le K}\chi_u(n)n^{-s}\right|^2
 \ll_{\mathcal K,\eps} Q^\eps K^{1-2\sigma}.
\end{equation}
The same estimate holds under any fixed local condition of positive limiting probability. No family member is removed.
\end{theorem}
\begin{proof}
Choose a smooth real $V$ equal to one on $(0,1]$ and zero on $[2,\infty)$, and put $P_K(s,u)=\sum_n\chi_u(n)n^{-s}V(n/K)$. Its Mellin transform $\widehat V(w)$ has a simple pole of residue one at $w=0$ and rapid decay on fixed vertical lines. Mellin inversion, a shift from $\Re(s+w)>1$ to $\Re(s+w)<0$, and the primitive functional equation give the exact identity
\begin{equation}\label{eq:flat-FE}
 L(s,\chi_u)-P_K(s,u)
 =(q/\pi)^{1/2-s}\sum_{n\ge1}\chi_u(n)n^{s-1}
                   W_s(\pi nK/q),
\end{equation}
where
\[
 W_s(y)=-\frac1{2\pi\ii}\int_{(-c)}\widehat V(w)
   \frac{\Gamma((2-s-w)/2)}{\Gamma((s+w+1)/2)}y^w\dd w.
\]
The only crossed pole in the first shift is $w=0$, since the characters are nonprincipal. The minus sign in the dual formula and the conductor factor are retained.

Moving the line defining $W_s$ left gives arbitrary polynomial decay for large $y$. Moving it to a sufficiently small positive line crosses only $w=0$, because the first numerator gamma pole has real part $2-\sigma>1$. Thus $W_s$ and its logarithmic derivative are bounded at zero. On a dyad $n\asymp V$, their supremum and variation are
\[
 \ll_J\min\{1,(VK/Q)^{-J}\},
\]
uniformly for $Q<u\le2Q$ and $s\in\mathcal K$.

Apply Lemma~\ref{lem:full-column-LS} to the dual with $\beta=1-\sigma$. The row norm of a dyad is at most
\[
 Q^{-h}(QV)^\eps(1+V/Q)^{1/2}V^h
                   \min\{1,(VK/Q)^{-J}\}.
\]
For $V\le Q/K$, dyadic Minkowski sums this to $O(Q^\eps K^{-h})$, since $h>0$. The dual unit has norm $O(Q^{-h})$ and also fits. Every dyad $V>Q/K$ is retained using $J>h+1$; this also pays those with $V>Q$. Finally, $P_K$ differs from the sharp initial segment only on $K<n\le2K$, where Lemma~\ref{lem:full-column-LS} with $\beta=\sigma$ gives the same row norm. Squaring proves \eqref{eq:L-tail}.
\end{proof}

A holomorphic submean argument gives the compact-supremum variant
\[
 \E_Q\sup_{s\in\mathcal K}
 \left|L(s,\chi_u)-\sum_{n\le K}\chi_u(n)n^{-s}\right|^2
 \ll_{\mathcal K,\eps}Q^\eps K^{1-2\sigma_0}.
\]
One uses a finite disc cover inside a slightly larger compact set with left edge $\sigma_0-\tau>1/2$, and absorbs $K^{2\tau}\le Q^{2\tau}$ after choosing preliminary losses. This asserts neither growing-height uniformity nor a critical-line endpoint.

\subsection{The limiting Euler object and a master residual comparison}
The random product
\[
 \mathcal L(s)=\prod_{p\text{ odd}}(1-X_pp^{-s})^{-1}
\]
exists as a locally holomorphic $L^2$-valued function on $\Re s>1/2$. In fact
\begin{equation}\label{eq:random-tail}
 \left\|\mathcal L(s)-\sum_{n\le K}X(n)n^{-s}\right\|_{L^2}
 \ll_{\mathcal K}K^{1/2-\sigma}\sqrt{\log(2K)}.
\end{equation}
To see the tail estimate, the absolute square kernel on one dyad is
\[
 \sum_{\substack{V<n,m\le2V\\nm=\square}}(nm)^{-\sigma}
 \ll V^{1-2\sigma}\log(2V).
\]
Write $n=ra^2,m=rb^2$ with common squarefree $r$, and use the resulting harmonic $r$-sum. Dyadic Minkowski gives \eqref{eq:random-tail}. Independently, the local means and second moments of the Euler factors are $1+O(p^{-2\sigma})$, so independence makes their products Cauchy in $L^2$ locally. Agreement with the series on $\Re s>1$ identifies the two continuations. Square aliases are essential to this argument.

Let $A_u(s)=\sum_{n\le X}c_n(s)\chi_u(n)$, with coefficients common to the rows, and let $A(s,X_\bullet)$ be its random counterpart. Define actual coefficient masses
\[
 B=\sum_{n\le X}\abs{c_n(s)},\qquad
 D=\sum_{n\le X}n^{1/4}\abs{c_n(s)}.
\]
For a bounded scalar $z=z(s)$ put
\[
 R_Q=\E_Q\abs{zL(s,\chi_u)A_u(s)-1}^2,\qquad
 R_\infty=\E\abs{z\mathcal L(s)A(s,X_\bullet)-1}^2.
\]

\begin{theorem}[Quantitative full-residual transfer]\label{thm:residual-transfer}
For $1\le K\le Q$,
\begin{align}
 \abs{\sqrt{R_Q}-\sqrt{R_\infty}}
 &\ll Q^\eps\abs z B K^{-h}
  +Q^{-1/4}\sqrt{\log(2KX)}(1+\abs z D K^\lambda),
 \label{eq:root-transfer}\\
 \abs{R_Q-R_\infty}
 &\ll Q^\eps(1+\abs z B)\abs z B K^{-h}
  +Q^{-1/2}\log(2KX)(1+\abs z D K^\lambda)^2.
 \label{eq:square-transfer}
\end{align}
The fixed-compact and fixed-local-conditioning conventions above apply.
\end{theorem}
\begin{proof}
Replacing the actual and random $L$-values by initial segments changes their residual norms by at most $O(Q^\eps\abs z B K^{-h})$, by \eqref{eq:L-tail} and \eqref{eq:random-tail}. Expand the truncated residual as a finite character polynomial $\sum_\ell b_\ell\chi_u(\ell)$, or its random analogue. Its support is at most $KX$, and
\[
 \sum_\ell\ell^{1/4}\abs{b_\ell}
 \le1+\abs z D\sum_{a\le K}a^{1/4-\sigma}
 \ll1+\abs z D K^\lambda.
\]
The complete square expansion and \eqref{eq:dyadic-character-mean} therefore bound the difference of truncated squared norms by
\[
 O\bigl(Q^{-1/2}\log(2KX)(1+\abs z D K^\lambda)^2\bigr).
\]
Use $\abs{\sqrt a-\sqrt b}\le\sqrt{\abs{a-b}}$ for nonnegative $a,b$ to obtain \eqref{eq:root-transfer}.

For \eqref{eq:square-transfer}, the actual full $L$ second moment is uniformly bounded on the compact set. This follows already from \eqref{eq:root-transfer} with $A=1,z=1,K=Q^{1/3}$ and the bounded random moment; it is not an input to that proof. Thus full residual norms are $O(1+\abs z B)$. The truncated $L$ norm is also $O(Q^{\eps'})$ by the tail estimate. Consequently the change of squared norms due to truncation is bounded by the first term of \eqref{eq:square-transfer}, after allocating the preliminary loss. Add the finite polynomial comparison. This argument never divides by $z$ or by any $L$-value.
\end{proof}

For the triangular inverse-Dedekind mollifier
\[
 \mathcal M_{u,X}(s)=\sum_{mn\le X}\mu(m)\mu(n)\chi_u(n)(mn)^{-s},
\]
the masses satisfy $B\ll X^{1-\sigma}\log(2X)$ and $D\ll X^{5/4-\sigma}$. Substitution in Theorem~\ref{thm:residual-transfer} gives the following strict ranges for $X\le Q^\eta$:
\begin{align}
 \eta&<\frac{\sigma_0-1/2}{5/2-2\sigma_0},&
 \delta_1&=\frac{\sigma_0-1/2-\eta(5/2-2\sigma_0)}3,
 \label{eq:root-range}\\
 \eta&<\frac{\sigma_0-1/2}{(5-4\sigma_0)(3/2-\sigma_0)},&
 \delta_2&=\frac{\sigma_0-1/2-\eta(5-4\sigma_0)(3/2-\sigma_0)}{4-2\sigma_0}.
 \label{eq:square-range}
\end{align}
The respective errors are $O(Q^{-\delta_1+\eps})$ for root residuals and $O(Q^{-\delta_2+\eps})$ for squared residuals. The choices are $K=Q^{(1-\eta)/3}$ and $K=Q^{(1-\eta)/(4-2\sigma_0)}$. At $\sigma_0=3/4$, these permit $\eta<1/4$ and $\eta<1/6$; at $5/6$, they permit $\eta<2/5$ and $\eta<3/10$. These fractions describe family approximation and mollifier lengths, not zero locations.

\subsection{Exact local optimization}
For $y\ge3$ let
\[
 E_y(s,X_\bullet)=\prod_{3\le p\le y}(1-X_pp^{-s}),\qquad
 \mathcal L=\mathcal L_{\le y}\mathcal T_y,
\]
and set $C_y=\E\mathcal T_y$, $H_y=\E\abs{\mathcal T_y}^2$. With $g_p=p/(p+1)$,
\begin{equation}\label{eq:Cy}
 C_y(s)=\prod_{p>y}\frac{1-(1-g_p)p^{-2s}}{1-p^{-2s}}.
\end{equation}
It is holomorphic and nonzero on $\Re s>1/2$.

\begin{theorem}[Optimal correction at a prescribed mean]\label{thm:variance}
Fix a scalar $m=m(s)$. Among all complex functions $N$ of the states $X_p$, $3\le p\le y$, subject to $\E[\mathcal L N]=m$, the residual satisfies
\begin{equation}\label{eq:variance-min}
 \min_N\E\abs{z\mathcal L N-1}^2
 =\abs{zm-1}^2+\abs{zm}^2\left(\frac{H_y}{\abs{C_y}^2}-1\right).
\end{equation}
The minimum is attained by $N_*=mE_y/C_y$. It is unique almost surely if $z\ne0$. If $z=0$, every residual is one and uniqueness is not asserted.
\end{theorem}
\begin{proof}
Put $A=\mathcal L_{\le y}N$. This is an invertible change of finite-state variables, and $\E A=m/C_y$. Independence of the tail gives exactly
\begin{align*}
 \E\abs{z\mathcal L N-1}^2
  =\abs{zm-1}^2
   +\abs{zm}^2\left(\frac{H_y}{\abs{C_y}^2}-1\right)
   +\abs z^2 H_y\E\abs{A-\E A}^2.
\end{align*}
The final term vanishes precisely when $A$ is constant, except for the stated degenerate case.
\end{proof}

This optimization is over the full finite local state space, not a selected linear subspace. Its remaining tail variance is explicit. For $x=p^{-s}$,
\[
 v_p(s)=\frac{g_p\abs x^2(1+(1-g_p)\abs x^2)}
                  {\abs{1-(1-g_p)x^2}^2},\qquad
 \frac{H_y}{\abs{C_y}^2}=\prod_{p>y}(1+v_p).
\]
Consequently
\begin{equation}\label{eq:variance-tail}
 0\le H_y/\abs{C_y}^2-1\ll_{\mathcal K}y^{1-2\sigma}.
\end{equation}
The last estimate only uses comparison of $\sum_{p>y}p^{-2\sigma}$ with the integer sum.

The correction also transfers to actual growing families. Put $P_y=\prod_{3\le p\le y}p$. Its coefficient masses are
\[
 B_y=\prod_{3\le p\le y}(1+p^{-\sigma}),\qquad
 D_y=\prod_{3\le p\le y}(1+p^{1/4-\sigma}).
\]
If $y=(\log Q)^\nu$ with $0<\nu<(5/4-\sigma_0)^{-1}$, then $B_y,D_y=Q^{o(1)}$ and $\log(2KP_y)=Q^{o(1)}$. This remains true even when $P_y$ exceeds every fixed power of $Q$; the master estimate uses the actual masses, not a small-support assertion. Taking $K=Q^{1/3}$ in \eqref{eq:root-transfer} yields
\begin{equation}\label{eq:actual-variance-removal}
 \E_Q\left|L(s,\chi_u)\frac{E_y(s,\chi_u)}{C_y(s)}-1\right|^2
 \ll y^{1-2\sigma}+Q^{-2(\sigma_0-1/2)/3+\eps}.
\end{equation}
For a fixed ordinary mollifier $M_Z(s)=\sum_{n\le Z}\mu(n)n^{-s}$, the complete common-factor residual therefore has limit
\begin{equation}\label{eq:remaining-scalar}
 \lim_{Q\to\infty}\E_Q\left|
 \zeta(s)L(s,\chi_u)M_Z(s)\frac{E_y(s,\chi_u)}{C_y(s)}-1
 \right|^2=\abs{\zeta(s)M_Z(s)-1}^2.
\end{equation}
The normalizer contains no inverse zeta factor. At an ordinary zero, both sides of \eqref{eq:remaining-scalar} equal one. Thus the auxiliary variance is removed, but the decisive strict-below-one scalar bound has not been proved.

\section{Two mollifiers and capped reciprocal moments}\label{sec:capped}

The following transfer, from \cite{R989}, is independent of the particular arithmetic family. It makes explicit why one polynomial need not perform both the control and leading-term tasks.

\begin{theorem}[Two-mollifier transfer]\label{thm:two-mollifier}
Let $L_i,P_i,Q_i$ be finite complex arrays, with bounded nonnegative weights $w_i$ of total mass $O(U)$. Fix positive integers $k,r$. Suppose, with arbitrarily small positive losses,
\begin{align*}
 \sum_iw_i(\abs{P_i}^{4k}+\abs{Q_i}^{4k})&\ll U^{1+\eps},\\
 \sum_iw_i\abs{L_iP_i-1}^{2r}&\ll U^{1-\delta+\eps},\\
 \sum_iw_i\abs{L_iQ_i-1}^{2}&\ll U^{1-\lambda+\eps},\\
 \sum_iw_i\abs{Q_i}^{2k}&=cU+O(U^{1-\tau+\eps}),
\end{align*}
where $\delta,\lambda,\tau>0$ are fixed. For $0<\theta<\delta$,
\begin{equation}\label{eq:two-mollifier}
 \sum_iw_i\min\{\abs{L_i}^{-2k},U^\theta\}
 =cU+O(U^{1-\eta+\eps}),\quad
 \eta=\min(\lambda/2,\tau,\delta/2,\theta,\delta-\theta).
\end{equation}
At $L_i=0$, the summand is defined to be the cap $U^\theta$.
\end{theorem}
\begin{proof}
Let $B=\{i:\abs{L_iP_i-1}>1/2\}$ and $G=B^c$. Markov gives $\sum_Bw_i\ll U^{1-\delta+\eps}$. Every zero of $L_i$ is in $B$. On $G$, $I_i=1/L_i$ exists and $\abs{I_i}\le2\abs{P_i}$. Put $F_i=L_iQ_i-1$. Since $I_i-Q_i=-I_iF_i$, the elementary power difference and Young inequalities give
\[
 \bigl|\abs{I_i}^{2k}-\abs{Q_i}^{2k}\bigr|
 \le C_k\abs{F_i}(\abs{P_i}^{2k}+\abs{Q_i}^{2k}).
\]
Cauchy--Schwarz bounds its weighted sum by $O(U^{1-\lambda/2+\eps})$. The missing $Q$ mass on $B$ is $O(U^{1-\delta/2+\eps})$, again by Cauchy--Schwarz.

The cap is not assumed inactive on $G$. Its complete loss is at most
\[
 U^{-\theta}\sum_Gw_i\abs{I_i}^{4k}\ll U^{1-\theta+\eps}.
\]
The capped contribution of $B$ is at most $U^\theta\sum_Bw_i\ll U^{1-\delta+\theta+\eps}$. Add these four errors to the short-polynomial moment asymptotic. There is no independence hypothesis between either polynomial and the family.
\end{proof}

\subsection{Source-qualified character-family application}
To state the scope of the arithmetic application, fix a field containing the $n$th roots of unity, $n\ge3$, admissible finite local data $S$, and one allowed fixed finite-order twist $\psi$. Rows are the actual $n$th-power-free exterior integral ideals, ordered by their full norms. Character values at nonunits are zero. Let
\[
 L_{\mathfrak a}(s)=L_F^S(s,\psi\chi_{\mathfrak a}),\qquad
 M_A(s,\mathfrak a)=\sum_{N\mathfrak d\le A}
 \frac{\mu_F(\mathfrak d)\psi(\mathfrak d)\chi_{\mathfrak a}(\mathfrak d)}{(N\mathfrak d)^s}.
\]
The application in \cite{R989} adopts the precise large-sieve, flat-cutoff functional-equation, full-residual, and arithmetic character-mean contracts from its pinned predecessors. In particular, its large-sieve kernel is
\begin{equation}\label{eq:n-sieve-contract}
 K_n(U,V)=U+V+U^{1-1/n}V^{2/n}+U^{2/n}V^{1-1/n},
\end{equation}
with power-free indices on both sides. Full columns $\mathfrak c=\mathfrak e\mathfrak f^n$ retain the exact mask
$\chi_{\mathfrak a}(\mathfrak c)=\chi_{\mathfrak a}(\mathfrak e)\ind_{(\mathfrak a,\mathfrak f)=1}$, without an extra condition $(\mathfrak e,\mathfrak f)=1$. The complete conjugate dual and its actual conductor are part of the functional-equation contract. These analytic contracts are not independently reconstructed in this manuscript; the next statement is explicitly conditional on them.

For a nonnegative smooth $W$ supported in $(1,2)$, put
\[
 c_W=\frac{\Res_{z=1}\zeta_F^S(z)}{\zeta_F^S(n)}\int_0^\infty W(x)\dd x,
 \qquad w_p=\frac{1-q^{-1}}{1-q^{-n}},\quad q=Np,
\]
and
\begin{equation}\label{eq:reciprocal-constant}
 C_k(s)=\prod_{p\notin S}\left[
 1-w_p+\frac{w_p}{n}\sum_{\omega^n=1}
               \abs{1-\psi(p)\omega q^{-s}}^{2k}\right].
\end{equation}
All neutral aliases for $k\ge n$ occur in this positive local average.

\begin{reported}[Capped fixed-order laws of \cite{R989}]\label{prop:all-capped}
Under the named source contracts, at $n=6$, $s=3/4+\ii t$, fixed $k\ge1$ and fixed compact $t$-range, define
\[
 e_k=\frac1{8k+2},\qquad
 \delta_k=\begin{cases}
 4/15,&k=1,\\1/4,&k=2,\\1/5,&3\le k\le5,\\
 1/k,&6\le k\le10,\\19/194,&11\le k\le15,\\3/(2k),&k\ge16.
 \end{cases}
\]
For every fixed $0<\theta<\delta_k$,
\begin{align}\label{eq:capped-application}
 \sum_{\mathfrak a}W(N\mathfrak a/U)
 \min\{\abs{L_{\mathfrak a}(s)}^{-2k},U^\theta\}
 =c_W C_k(s)U
 +O\left(U^{1-\min(e_k,\theta,\delta_k-\theta)+\eps}\right).
\end{align}
More generally, on every fixed line $1/2<\sigma<1$, put $d=2\sigma-1$, $a=\min(d/2,1/(4k))$, and $\Delta=da$. The corresponding source-qualified cap and error are $U^{3\Delta/4}$ and $O(U^{1-\Delta/4+\eps})$.
\end{reported}

The arithmetic leading term is obtained by expanding a short polynomial before averaging. Nonprincipal tuples cost
$U^{1/2+\eps}B^{k(5/2-2\sigma)+\eps}$, while the omitted neutral-tuple tail costs $UB^{-d/2}$. The long polynomial is chosen using the predecessor's complete residual estimate. Theorem~\ref{thm:two-mollifier} then supplies the displayed transfer, including the cap loss on the good set. This explanation identifies the operation contributed by \cite{R989}; it is not a substitute for its deep source contracts.

At $k=1$, the two useful choices are cap $U^{1/6}$ with error $U^{9/10+\eps}$ and cap $U^{1/4}$ with error $U^{59/60+\eps}$. Every statement is for fixed order. The constants are not controlled for an unbounded-order limit, and the cap exponent at large $k$ is not uniformly superior to every earlier fixed cap. At a zero, the cap remains compatible with the family error. This is an all-fixed-moments result of a capped kind, not an uncapped moment theorem or a zero-free half-plane.

\section{Keeping the signs in M\"obius sums}\label{sec:native}

The family methods leave an exceptional-member problem. A different approach works directly with the ordinary M\"obius function. The starting point is an exact identity: knowing $\mu(n)$ up to $Y$ determines a quadratic expression that agrees with $\mu(n)$ up to, but not including, $(Y+1)^2$.

This is useful only if its signs are retained. Bounding every summand by its absolute value gives an estimate too large for the intended application. The work in \cite{R994,R1003} separates contributions that can already be bounded from the remaining signed sum. In the repository this is called the \emph{native} approach: the coefficients are the original M\"obius coefficients, not arbitrary bounded substitutes. We retain that meaning here.

\subsection{An exact identity from one scale to its square}
Fix an integer $Y\ge2$, put $b=Y+1$, and define
\[
 g_Y(n)=\mu(n)\ind_{n\le Y},\qquad z_Y=g_Y*g_Y,\qquad
 G_Y(s)=\sum_{n\le Y}\mu(n)n^{-s}.
\]
The complete reduced-denominator coefficient is
\begin{equation}\label{eq:CY}
 C_Y(q)=\sum_{\substack{d\le Y^2\\q\mid d}}\frac{z_Y(d)}d.
\end{equation}
Products above the eventual output variable are included in this definition. Put
\[
 R(x)=h_{\lfloor x\rfloor}-\log x-\gamma,\qquad h_0=0,
\]
and, for $q\ge2$,
\begin{equation}\label{eq:Sq}
 S_q(x)=\sum_{n\le x}\frac{c_q(n)}n+\Lambda(q)
       =\sum_{d\mid q}\mu(q/d)R(x/d),
\end{equation}
where $c_q$ is the Ramanujan sum. The equality follows from
$c_q(n)=\sum_{d\mid(q,n)}d\mu(q/d)$ and
$\sum_{d\mid q}\mu(q/d)\log d=\Lambda(q)$. In particular the logarithmic and constant terms cancel, so $S_q$ is constant on complete integer cells.

Let
\[
 M_Y=m(Y),\quad L_Y=\sum_{n\le Y}\frac{\mu(n)\log n}{n},\quad
 B_Y=1+L_Y-(\log Y+\gamma/2)M_Y,
\]
and $A_Y(x)=2M_YB_Y+M_Y^2\log(Y^2/x)$. Then
\begin{equation}\label{eq:Newton-exact}
 m(k)=A_Y(k)-M_Y^2R(k)-\sum_{q=2}^{Y^2}C_Y(q)S_q(k),
 \qquad b\le k<b^2.
\end{equation}
Indeed the Dirichlet-convolution identity
\[
 \mu-(2g_Y-\mathbf1*g_Y*g_Y)
 =\mu*(\delta-\mathbf1*g_Y)* (\delta-\mathbf1*g_Y)
\]
has right side supported on $n\ge b^2$, since each factor $\delta-\mathbf1*g_Y$ vanishes below $b$. Here $\mathbf1(n)=1$ and $\delta$ is the convolution unit. Weighted summation gives
$m(k)=2M_Y-\sum_d z_Y(d)h_{\lfloor k/d\rfloor}/d$. Grouping the reduced rational frequencies gives \eqref{eq:Newton-exact}. The first possible coefficient defect at $b^2$ is $\mu(b)^2$, so the endpoint is strict.

\subsection{Factoring the reduced-denominator coefficients}
For $d\ge1$ put $m_d(x)=\sum_{n\le x,(n,d)=1}\mu(n)/n$, with value zero for $x<1$. We use Tao's elementary bound \cite{Tao10}
\begin{equation}\label{eq:masked-m}
 \abs{m_d(x)}\le1.
\end{equation}
For completeness, at integral $x$ let $A$ count integers at most $x$ coprime to $d$, and let $B$ count those whose prime factors all divide $d$. Divisor inversion gives
$B=\sum_{n\le x,(n,d)=1}\mu(n)\lfloor x/n\rfloor$. Since $\{x/n\}\le1-1/n$,
\[
 \abs{xm_d(x)}\le B+A-\sum_{n\le x,(n,d)=1}\frac1n\le x.
\]
The two counted sets intersect only at one. Real $x$ follows by taking its integer part.

\begin{proposition}[Complete denominator coefficients]\label{prop:denominator}
If $q$ has a prime cube factor, $C_Y(q)=0$. Otherwise write uniquely $q=r^2s$, with $r,s$ coprime and squarefree. Then
\begin{equation}\label{eq:coefficient-factorization}
 C_Y(r^2s)=\frac{\mu(s)}{r^2s}
 \sum_{\substack{uvw=s\\u,v,w\text{ pairwise coprime}}}
 \frac{m_{ruw}(Y/(ruw))\,m_{rvw}(Y/(rvw))}{w}.
\end{equation}
Consequently
\begin{equation}\label{eq:coefficient-bounds}
 \abs{qC_Y(q)}\le\prod_{p\mid s}(2+1/p),\qquad
 \sum_{q\le Q}q^2\abs{C_Y(q)}^2\le32Q(1+\log Q)^3.
\end{equation}
\end{proposition}
\begin{proof}
Both old inputs are squarefree, so their product is cubefree. For squarefree $d$ define
$T(d)=\sum_{a\le Y,d\mid a}\mu(a)/a=\mu(d)m_d(Y/d)/d$. Every prime of $r$ must divide both inputs. For $p\mid s$ expand
\[
 \ind_{p\mid ac}=\ind_{p\mid a}+\ind_{p\mid c}
                  -\ind_{p\mid a}\ind_{p\mid c}.
\]
Assigning these three possibilities to $u,v,w$ gives
$C_Y(r^2s)=\sum_{uvw=s}(-1)^{\omega(w)}T(ruw)T(rvw)$. Each prefactor has sign $\mu(s)$ and denominator $r^2sw$, proving \eqref{eq:coefficient-factorization}. Bound both masked sums by \eqref{eq:masked-m} to obtain the first inequality.

For the second, let $a(q)$ be the square of the first majorant, zero off cubefree integers. Locally, $a(p)=(2+1/p)^2$ and $a(p^2)=1$. Thus
\[
 a(q)\le d_4(q)\prod_{p\mid q}\left(1+\frac1p+\frac1{4p^2}\right).
\]
Expand the positive squarefree divisor product and use
$d_4(ab)\le d_4(a)d_4(b)$ and $\sum_{n\le x}d_4(n)\le x(1+\log x)^3$. The remaining Euler constant is
\[
 \prod_p(1+4p^{-2}+p^{-3})<32.
\]
Indeed its logarithm is at most $4\sum_{n\ge2}n^{-2}+\sum_{n\ge2}n^{-3}\le13/4$, using integral upper bounds, and $e^{13/4}<32$. This pays the full infinite prime tail.
\end{proof}

The rectangular factorization \eqref{eq:coefficient-factorization} relies on all products in \eqref{eq:CY}. An additional cutoff $d\le k$ inside that sum would destroy the factorization, and is not part of this theorem.

\subsection{The denominator ranges that can be controlled}
The low-denominator estimate of \cite[native proof, Lemma 3]{R994} is
\begin{equation}\label{eq:low-denominator}
 \sum_{K\le k<2K}\left|\sum_{2\le q\le Q}C_Y(q)S_q(k)\right|^2
 \le64(1+\log Q)^4\left(\frac QK+\frac{Q^3}{K^2}\right),
 \quad 1\le Q\le K.
\end{equation}
Here the square is taken after the full sum. To indicate its proof and the normalization, for $0<\vartheta<1$ write
\[
 Z_\vartheta(k)=\sum_{n\le k}\frac{e^{2\pi\ii n\vartheta}}n
                      +\log(1-e^{2\pi\ii\vartheta}).
\]
Abel summation gives
\[
 Z_\vartheta(k)=
 \frac{e^{2\pi\ii k\vartheta}}{(k+1)(1-e^{-2\pi\ii\vartheta})}
 +E_\vartheta(k),\qquad
 \abs{E_\vartheta(k)}\le\frac{2}{k^2\abs{1-e^{2\pi\ii\vartheta}}^2}.
\]
Summing primitive fractions of denominator $q$ recovers $S_q$. The distinct fractions with denominator at most $Q$ are $Q^{-2}$-separated. Their complete Gram row bound is $K+2Q^2(1+\log Q)$, while their squared coefficient mass is at most $\tfrac14\sum_{q\le Q}q^2\abs{C_Y(q)}^2$. Equation~\eqref{eq:coefficient-bounds} and the entire Abel remainder yield \eqref{eq:low-denominator}. In particular neither same-frequency aliases nor cross-denominator terms are dropped.

Choose $1/2<\rho<2/3$, put $\alpha=2-3\rho$, and partition $[b,b^2)$ into dyadic intervals $I_j=[X_j,U_j)$ with $X_j=2^jb$, $U_j=\min(2X_j,b^2)$. On the integer cells in $I_j$ set
\[
 V_{Y,\rho}(k)=-\sum_{X_j^\rho<q\le Y^2}C_Y(q)S_q(k).
\]
Summing \eqref{eq:low-denominator} over these complete intervals gives
\begin{equation}\label{eq:native-error}
 \norm{m-A_Y+M_Y^2R-V_{Y,\rho}}_2^2
 \le\frac{2048}{1-2^{-\alpha}}b^{-\alpha}(1+\log b)^4.
\end{equation}
The norm is over $b,\ldots,b^2-1$. A shortened last interval is paid by its full positive dyadic bound.

The large-denominator sector can also be reduced without assuming cancellation between different denominators.

\begin{theorem}[Complete small-core bound]\label{thm:small-core}
Let $\omega_q(x)$ be arbitrary measurable functions with $\abs{\omega_q(x)}\le1$. Uniformly for integers $Y\ge2$ and $S\ge1$,
\begin{align}
 \left\|\sum_{\substack{2\le r\le Y\\r\text{ squarefree}}}
       \omega_{r^2}C_Y(r^2)S_{r^2}\right\|_{L^2(0,\infty)}
 &\le3\sqrt3\,h_Y,\label{eq:square-sector}\\
 \left\|\sum_{\substack{2\le q\le Y^2\\q\text{ cubefree}\\
                    \mathrm{core}(q)\le S}}
       \omega_q C_Y(q)S_q\right\|_{L^2(0,\infty)}
 &\le12h_Y\sqrt S\,h_S^3.
 \label{eq:core-sector}
\end{align}
Here $\mathrm{core}(r^2s)=s$ for coprime squarefree $r,s$.
\end{theorem}
\begin{proof}
One has $\int_0^\infty R(x)^2\dd x\le3$: on $x\ge1$ use $\abs{R(x)}\le1/x$, and on $(0,1)$ integrate $(-\log x-\gamma)^2$. Dilation in \eqref{eq:Sq} gives
\[
 \norm{S_q}_2\le\sqrt{3q}\,\beta(q),\qquad
 \beta(q)=\prod_{p\mid q}(1+p^{-1/2}).
\]
For a square denominator, \eqref{eq:coefficient-factorization} gives
$C_Y(r^2)=m_r(Y/r)^2/r^2\ge0$. Expanding $\beta(r)$ on squarefree $r$ yields
\[
 \sum_{\substack{r\le Y\\r\text{ squarefree}}}\frac{\beta(r)}r
 \le\sum_{e\le Y}e^{-3/2}h_{\lfloor Y/e\rfloor}
 \le3h_Y.
\]
The triangle inequality proves \eqref{eq:square-sector}, and multiplication by $\omega_q$ contracts every individual norm.

For $q=r^2s$, put $a(s)=\prod_{p\mid s}(2+1/p)$. The corresponding norm is at most
$\sqrt3\,a(s)\beta(s)\beta(r)/(r\sqrt s)$. Sum over $r$ first. For squarefree $s$, one has $a(s)\beta(s)\le(9/8)d_4(s)$: the local factor is below four at every $p\ge3$, and below $9/2$ at $p=2$. Partial summation and the divisor bound give
\[
 \sum_{s\le S}\frac{d_4(s)}{\sqrt s}\le2\sqrt S\,h_S^3.
\]
The resulting constant is $27\sqrt3/4<12$. No diagonalization of the coherent sum was used.
\end{proof}

At the literal reciprocal sign crossings
\[
 \mathcal X=\{Y\ge2:\mu(Y)\ne0,\ m(Y-1)m(Y)\le0\},
\]
one has $\abs{M_Y}\le1/Y$. This set is unbounded by the classical Landau argument applied to
$\int_1^\infty m(x)x^{-z-1}\dd x=1/[z\zeta(z+1)]$: eventual one-sidedness would force a real singularity, whereas this expression is analytic at every real $z>-1$ and has nonreal poles there. No bound on gaps between crossings is used.

The elementary identity
\[
 B_Y=(\gamma/2)M_Y+\sum_{n\le Y}\frac{\mu(n)}nR(Y/n)
\]
gives $\abs{B_Y}\le1+(\gamma/2)\abs{M_Y}$. Moreover $B_Y\to0$ by the prime number theorem: split the last sum at $\eta Y$, bound the first part by $\eta$, and apply partial summation with $M(x)=o(x)$ to the bounded-variation function $R(1/u)/u$ on $[\eta,1]$. Combining this with \eqref{eq:native-error} proves, on $\mathcal X$,
\begin{equation}\label{eq:crossing-comparison}
 \left|\sqrt{F_{b^2-1}-F_Y}-\norm{V_{Y,\rho}}_2\right|
 \le\eta_Y,\qquad \eta_Y\longrightarrow0.
\end{equation}
If $V^{\mathrm{rough}}_{Y,\rho,S}$ retains only cores greater than $S$, Theorem~\ref{thm:small-core} permits the actual annular masks in one norm and gives
\begin{equation}\label{eq:rough-comparison}
 \left|\sqrt{F_{b^2-1}-F_Y}-\norm{V^{\mathrm{rough}}_{Y,\rho,S}}_2\right|
 \le\eta_Y+12h_Y\sqrt S\,h_S^3.
\end{equation}
There is no additional count of annuli. For $\log S=o(\log Y)$ the squared removed cost is subpower, not necessarily vanishing.

\subsection{A bound for the entire high-frequency tail}
The additional result of \cite[arithmetic frequency-tail proof]{R1003} has a short direct derivation. Let
\[
 Z_Y(x)=\sum_{d\le Y^2}\frac{z_Y(d)}dR(x/d),\qquad
 a_Y(u)=e^{u/2}Z_Y(e^u).
\]
With Fourier convention $\widehat a(t)=\int_\R a(u)e^{-\ii tu}\dd u$, the exact Mellin identity is
\begin{equation}\label{eq:native-Fourier}
 \widehat a_Y(t)=\Phi_Y(t)
 :=\frac{\zeta(1/2+\ii t)G_Y(1/2+\ii t)^2}{-1/2+\ii t}.
\end{equation}
It follows by applying the Mellin transform to $R$ and then to the finite dilation sum. The original products and both square factors of $G_Y$ are retained.

\begin{proposition}[High-frequency removal]\label{prop:frequency-tail}
Suppose $\abs{\zeta(1/2+\ii t)}\ll_\eps(1+\abs t)^{\vartheta+\eps}$ for a fixed $0\le\vartheta<1/2$. Then, for $T\ge1$,
\begin{equation}\label{eq:frequency-tail-general}
 \int_{\abs t>T}\abs{\Phi_Y(t)}^2\dd t
 \ll_\eps (YT)^\eps T^{-2+2\vartheta}(T+Y^2).
\end{equation}
With the published exponent $\vartheta=13/84$ of Bourgain \cite{Bourgain17}, this becomes
\[
 \int_{\abs t>T}\abs{\Phi_Y(t)}^2\dd t
 \ll_\eps(YT)^\eps T^{-71/42}(T+Y^2),
\]
and the entire tail at $T=Y^{84/71}$ has squared norm $O_\eps(Y^\eps)$.
\end{proposition}
\begin{proof}
Expanding $G_Y^j$ gives coefficients bounded by $d_j(n)$ on $n\le Y^j$. The time Gram matrix on $[T,2T]$ has row bound $O(T+Y^j\log(2Y))$. Also $d_j(n)^2\le d_{j^2}(n)$, as seen locally by counting nonnegative $j$ by $j$ matrices with prescribed row and column sums. Thus
\[
 \int_T^{2T}\abs{G_Y(1/2+\ii t)}^{2j}\dd t
 \ll_j(T+Y^j)(\log(2Y))^{j^2+1}.
\]
Apply this with $j=2$ and bound the zeta factor and denominator pointwise. Each dyad contributes
$O_\eps((YU)^\eps U^{-2+2\vartheta}(U+Y^2))$. Both powers sum geometrically for a sufficiently small preliminary loss, proving \eqref{eq:frequency-tail-general}.
\end{proof}

Frequency projection precedes physical restriction. If $a_{Y,\le T}$ is the inverse Fourier transform of $\ind_{\abs t\le T}\Phi_Y(t)$ and $W_Y=\ind_{[\log b,2\log b)}$, Plancherel gives
\[
 \norm{W_Y(a_Y-a_{Y,\le T})}_2^2
 \le\frac1{2\pi}\int_{\abs t>T}\abs{\Phi_Y(t)}^2\dd t.
\]
At the crossings, the complete strict-square-step identity therefore implies
\[
 \left|\sqrt{F_{b^2-1}-F_Y}-\norm{W_Ya_{Y,\le T}}_2\right|
 \ll1+\left(\int_{\abs t>T}\abs{\Phi_Y(t)}^2\dd t\right)^{1/2}.
\]
After multiplication by $W_Y$, the low-frequency component is not band-limited; its Fourier transform is a convolution. The bound above handles all mixed terms through a root-norm comparison.

\subsection{What an intermediate contraction would require}
The energy-adaptive cutoff of \cite{R1002} permits spending a specified error allowance $G_Y>0$ on a larger core sector. Set
\[
 S_Y=\min\left(Y^2,\left\lfloor
       \frac{G_Y}{2304h_Y^2h_{Y^2}^6}\right\rfloor\right).
\]
The value zero removes no core. Equation~\eqref{eq:core-sector} gives exactly
$\norm{V_{Y,\rho}-V^{\mathrm{rough}}_{Y,\rho,S_Y}}_2\le\sqrt{G_Y}/4$.
For $G_Y=Y^a(1+F_Y)^p$, this leaves an energy-dependent high-core target. It does not estimate that target.

A separate, elementary conditional statement explains the numerical goal. If a nonnegative nondecreasing sequence $F$ has a polynomial upper bound and satisfies, for \emph{every sufficiently large integer} $Y$,
\begin{equation}\label{eq:recurrence}
 1+F_{(Y+1)^2-1}\le C Y^a(1+F_Y)^p,
 \qquad a\ge0,\quad 1\le p<2,
\end{equation}
then $F_X\ll_\eps X^{a/(2-p)+\eps}$. Indeed an input exponent $\kappa$ is replaced by $(a+p\kappa)/2$; after $j$ fixed iterations it is
\[
 \frac a{2-p}+\left(\frac p2\right)^j
                   \left(\kappa_0-\frac a{2-p}\right).
\]
For each desired loss choose a sufficiently large fixed $j$, then the earlier losses. For the actual $m$, dyadic Cauchy--Schwarz makes its Mellin transform holomorphic on $\Re z>(\kappa-1)/2$ when $F_X\ll X^{\kappa+\eps}$. This excludes zeta zeros to the right of $(1+\kappa)/2$.

Thus $(a,p)=(1/3,3/2)$ would give $5/6$, and $(1/4,3/2)$ would give $3/4$. Neither recurrence is established here. In particular, \eqref{eq:crossing-comparison} on a sparse cofinal set is not itself the all-$Y$ premise \eqref{eq:recurrence}. The separate all-prefix growth theorem invoked in \cite{R994} must be checked at its stated hypotheses before using a weaker cofinal formulation.

The low-frequency counterexample in \cite[\texttt{LOW\_FREQUENCY\_ATTACK.md}]{R1003} further distinguishes the arithmetic task from a generic operator inequality. It constructs bounded squarefree-supported coefficient arrays with first coefficient one and both smooth balance coordinates zero, but with old energy of order $Y^\kappa$ and output energy of order $Y^{2\kappa}$. Those arrays do not obey M\"obius inversion. The result rules out obtaining the advertised contraction from that larger coefficient class alone; it is not a counterexample for the actual M\"obius coefficients.

\section{Approximating the xi function by functions with real zeros}\label{sec:theta}

The Riemann hypothesis can also be approached as an approximation problem. Write $\Xi(z)=\xi(1/2+\ii z)$. Its zeros are all real exactly when the nontrivial zeros of $\zeta$ lie on the critical line. A sequence of functions with real zeros would therefore be useful if it converged to $\Xi$ in a sufficiently strong sense.

The finite result is concrete. The constructions in \cite{R1003} match the theta expansion through Taylor orders $22$, $24$, and $26$ using a Gaussian times an even polynomial with real zeros. This goes beyond the order-$20$ Ising construction of \cite{R994}, but uses a different class of functions: a positive probability-density representation is no longer required. The change enlarges the available constructions; it does not identify them with ferromagnetic models.

The main distinction in this section is between \emph{matching a fixed number of coefficients} and \emph{matching arbitrarily many coefficients with a convergence theorem}. We have finite certificates for the former and a sufficient criterion for the latter. The all-order existence statement remains open.

\subsection{The theta integral and its normalization}
For $t\ge0$ put $q_n=\pi n^2$ and
\begin{equation}\label{eq:theta-density}
 \Phi(t)=\sum_{n\ge1}
       (4q_n^2e^{9t/2}-6q_ne^{5t/2})e^{-q_ne^{2t}},
\end{equation}
and extend evenly. Every summand is positive on $t\ge0$. Normalizing the density gives a random variable $T$ with
\[
 F_\theta(z)=\E e^{\ii zT}
 =\frac{\xi(1/2+\ii z)}{\xi(1/2)},\qquad
 v_\theta=\E T^2=0.046209986230837941577867620860678\ldots.
\]
The decimal is descriptive; the source uses directed intervals for the full integral. Define
\[
 H(z)=F_\theta(z/\sqrt{v_\theta}),\qquad H(0)=1,\quad H''(0)=-1.
\]
All matching equations in this section concern $H$. A physical approximant to $F_\theta(z)$ is obtained by replacing $z$ by $\sqrt{v_\theta}z$. In particular, a band of width $b$ in the standardized variable becomes a band of width $b/\sqrt{v_\theta}$ in the physical variable.

Let $\kappa_{2r}$ be the even cumulants of $T/\sqrt{v_\theta}$, and set
\begin{equation}\label{eq:Qr}
 Q_r=\frac{(-1)^{r+1}r\kappa_{2r}}{(2r)!},\qquad r\ge2.
\end{equation}
These numbers are defined by the theta moments, without a zero sum or an assumed real-zero product representation. On a disc about zero,
\begin{equation}\label{eq:logH}
 \log H(z)=-\frac{z^2}{2}-\sum_{r\ge2}\frac{Q_r}{r}z^{2r}.
\end{equation}

\subsection{A finite power-sum construction}
\begin{proposition}\label{prop:finite-P}
Suppose $b_1,\ldots,b_N>0$, $a=1/2-\sum_jb_j>0$, and
\begin{equation}\label{eq:power-sum-target}
 \sum_{j=1}^N b_j^r=Q_r,\qquad 2\le r\le M.
\end{equation}
Then
\[
 P(z)=e^{-az^2}\prod_{j=1}^N(1-b_jz^2)
\]
is even and entire, has only real zeros, and satisfies
\begin{equation}\label{eq:P-envelope}
 P(z)-H(z)=O(z^{2M+2}),\qquad
 \abs{P(z)}\le e^{\abs z^2/2}\quad(z\in\C).
\end{equation}
\end{proposition}
\begin{proof}
The nonzero exponential contributes no zeros, and each polynomial factor has roots $\pm b_j^{-1/2}$. Near zero,
\[
 \log P(z)=-\frac{z^2}{2}
       -\sum_{r\ge2}\frac{z^{2r}}r\sum_jb_j^r.
\]
Comparison with \eqref{eq:logH} proves Taylor matching. Finally,
$\abs{e^{-az^2}}\le e^{a\abs z^2}$ and
$\abs{1-b_jz^2}\le e^{b_j\abs z^2}$. Multiply and use $a+\sum b_j=1/2$.
\end{proof}

\begin{certificate}[Frozen source through order 26]\label{cert:theta26}
The directed data and proofs at the immutable source \cite{R1003} provide five solutions of \eqref{eq:power-sum-target}, with strict $a>0$ and $b_j>0$, as listed in Table~\ref{tab:theta-certificates}. Each active-coordinate cube has radius $10^{-50}$. The two incremental endpoints have twenty equal newest nodes, all fixed during their individual certificate calculations.
\end{certificate}

\begin{table}[ht]
\centering
\caption{Recorded directed theta certificates. The bounds are rounded upward from the source receipts; $\beta$ and $L$ are defined in \eqref{eq:Banach-bounds}. The Taylor order is not the polynomial degree or a verified zero height.}\label{tab:theta-certificates}
\small
\begin{tabular}{lrrrrr}
\toprule
Construction & Order $2M$ & Factors $N$ & Active & $\beta<$ & $L<$\\
\midrule
Direct &22&64&10&$1.155\cdot10^{-92}$&$2.320\cdot10^{-38}$\\
Direct &24&128&11&$1.518\cdot10^{-90}$&$5.857\cdot10^{-37}$\\
Direct &26&128&12&$1.194\cdot10^{-87}$&$8.859\cdot10^{-35}$\\
Incremental endpoint &24&84&11&$4.499\cdot10^{-90}$&$1.749\cdot10^{-36}$\\
Incremental endpoint &26&104&12&$1.231\cdot10^{-87}$&$9.236\cdot10^{-35}$\\
\bottomrule
\end{tabular}
\end{table}

Here is the precise finite-to-infinite argument used by those certificates. The parameter files specify exact rational centers $x_0$, active indices, and a rational preconditioner $R$. All inactive nodes remain fixed. For the vector
\[
 f(x)=\left(\sum_jb_j^r-Q_r\right)_{r=2}^M,
 \qquad \frac{\partial f_r}{\partial b_j}=r b_j^{r-1},
\]
the interval calculation covers every term of each finite sum and every Jacobian entry on the full cube. It verifies positivity throughout and gives
\begin{equation}\label{eq:Banach-bounds}
 \norm{Rf(x_0)}_\infty\le\beta,\qquad
 \sup_x\norm{I-RDf(x)}_\infty\le L<1,\qquad
 \beta+L\rho<\rho,
 \quad \rho=10^{-50}.
\end{equation}
The map $x\mapsto x-Rf(x)$ is a strict contraction into the cube. Moreover $RDf(x_0)$ is invertible by the norm bound, so $R$ is invertible. Its fixed point therefore solves $f(x)=0$. This proves existence of a true solution, not equality of a displayed decimal center with that solution.

The theta directory at the pin in Appendix~\ref{app:provenance} contains the rational centers in \nolinkurl{gaussian_parameters.json} and \nolinkurl{incremental_parameters.json}. The accepting routines are \nolinkurl{check_gaussian.py} and \nolinkurl{check_incremental.py}. The complete moment source has SHA-256
\begin{center}\footnotesize\texttt{c171903e4f60c189a929ccfcce36f7db034508e9ca10c85ed766338625558dcf}.
\end{center}
The mathematical acceptance requires the source bounds described next; a matching hash alone is not a proof of them. This manuscript records the frozen certificates and their proof mechanism. Their full order-26 source reconstruction was not newly executed during the editorial compilation.

\subsection{Bounding the errors in the theta integrals}
The producer described in \cite[\texttt{theta/SOURCE26.md}]{R1003} uses an authenticated interval backend from \cite{R994}, configured for 512 fractional bits and Taylor degree 220. It explicitly retains indices $n=1,\ldots,11$ and integrates to $U_n=c_n/16$, where
\[
 (c_1,\ldots,c_{11})=(40,29,23,18,14,12,9,7,5,3,2).
\]
There are 162 cells of width $1/16$, each with half-width $1/32$. The full density summand is expanded on a Cauchy circle of radius $1/8$, then multiplied by $t^j$ and integrated exactly for each even $j\le26$. On these circles $q_n<484$, $\Re t<21/8$, and $\abs{\Im t}\le1/8$, so the summand is bounded in modulus by $10^{16}$. The complete full-line Taylor integration error is at most
\begin{equation}\label{eq:theta-Taylor-error}
 E_j=\frac{32\cdot10^{16}(5/2)^j}{4^{221}(3/4)}.
\end{equation}
This uses the ratio $1/4$ of the cell half-width to the Cauchy radius and an upper bound 32 for the total full-line interval length.

For a retained index let $q=\pi n^2$, $E=e^{2U}$, and $\lambda=2qE-9/2>0$. The full omitted time interval satisfies
\begin{equation}\label{eq:theta-time-error}
 2\int_U^\infty t^j\Phi_n(t)\dd t
 \le8q^2e^{9U/2-qE}
       \sum_{k=0}^j\binom jk\frac{U^{j-k}k!}{\lambda^{k+1}}.
\end{equation}
Indeed $e^{2x}\ge1+2x$ for $x\ge0$ supplies an exponential tail after writing $t=U+x$, and dropping the negative prefactor term gives a positive majorant. For all omitted indices $n\ge12$, the complete bound is
\begin{equation}\label{eq:theta-index-error}
 \sum_{n\ge12}2\int_0^\infty t^j\Phi_n(t)\dd t
 \le\frac{16\pi^2 12^4e^{-144\pi}j!}{(288\pi-9/2)^{j+1}}.
\end{equation}
The consecutive numerator ratio $(1+1/n)^4e^{-\pi(2n+1)}$ is below $1/2$; the denominator increases. This pays the entire index tail.

After the normalization interval is proved positive, directed division gives the raw normalized moments and their variance-one versions. Cumulants are formed by the exact finite recursion
\[
 \kappa_n=\mu_n-\sum_{j=1}^{n-1}\binom{n-1}{j-1}\kappa_j\mu_{n-j}.
\]
Every Taylor, time, and index remainder enters the target intervals used in \eqref{eq:Banach-bounds}. Public zero-table entries were used only to initialize exploratory numerical searches; the accepting checker does not load a zero table. The successful numerical chain $64/22\to84/24\to104/26$ is not a certified connecting path: each endpoint has its own certificate, and no unbounded induction follows from two successful extensions.

\subsection{What would suffice at unbounded orders}
The next theorem is the all-order component of \cite[\texttt{WEIGHTED\_CANONICAL.md}]{R1003}. It removes the need for a positive Gaussian reserve, a positive inverse Fourier density, integer multiplicities, or a common support bound. Its unbounded feasibility hypothesis is not proved by the five finite certificates.

\begin{theorem}[Cofinal positive weighted feasibility]\label{thm:weighted-canonical}
Suppose that for an unbounded sequence of integers $M$ there are finitely many positive real nodes $b_{M,j}$ and positive real weights $\lambda_{M,j}$ with
\begin{equation}\label{eq:weighted-moments}
 \sum_j\lambda_{M,j}b_{M,j}^{\,r}=Q_r,\qquad 2\le r\le M.
\end{equation}
Then every zero of $H$ is real, and RH follows. No uniform bound on node count, support, total weight, or $\sum_j\lambda_{M,j}b_{M,j}$ is assumed.
\end{theorem}
\begin{proof}
On the open upper half-plane, $1-bz^2$ avoids the nonpositive real axis for every $b>0$. The principal branch $\ell_b(z)=\operatorname{Log}(1-bz^2)$ is holomorphic there. Define
\[
 P_M(z)=\exp\left[-\frac{z^2}{2}
       +\sum_j\lambda_{M,j}\{\ell_{b_{M,j}}(z)+b_{M,j}z^2\}\right].
\]
It is holomorphic and nonvanishing on this half-plane. For noninteger weights it need not be entire; its branch points are on the real boundary.

The elementary canonical-factor inequality
\begin{equation}\label{eq:canonical-inequality}
 \abs{(1-w)e^w}\le\exp(\abs w^2/2),\qquad w\in\C,
\end{equation}
follows from $\abs{1-w}^2=1-2\Re w+\abs w^2$ and $\log(1+u)\le u$; $w=1$ is immediate. Positive weights and the $r=2$ equation give
\begin{equation}\label{eq:quartic-envelope}
 \abs{P_M(z)}\le\exp\left(\frac{\abs z^2}{2}
                              +\frac{Q_2\abs z^4}{2}\right).
\end{equation}
This establishes normality without a support restriction.

One must still identify the limit without assuming a common Taylor disc: large nodes could bring branch points arbitrarily close to zero. For $x\ge0$ the exact logarithm remainder is
\[
 \log(1+x)-x=\sum_{r=2}^{M-1}\frac{(-1)^{r+1}x^r}{r}+R_M(x),
 \qquad \abs{R_M(x)}\le x^M/M.
\]
The estimate is valid for all $x$, not just $x<1$. At $z=\ii t$, the real logarithm therefore satisfies
\[
 \left|\log P_M(\ii t)-\left(
 \frac{t^2}{2}+\sum_{r=2}^{M-1}\frac{(-1)^{r+1}Q_rt^{2r}}r\right)\right|
 \le\frac{Q_Mt^{2M}}M.
\]
Choose $d>0$ inside a disc on which $\log H$ is analytic, and let $C=\max_{\abs z=d}\abs{\log H(z)}$. Cauchy's bound gives $\abs{Q_r}/r\le Cd^{-2r}$. For every fixed $0<t<d$, both the displayed remainder and the tail of \eqref{eq:logH} tend to zero. Thus $P_M(\ii t)\to H(\ii t)$ on an interval inside the upper half-plane.

By \eqref{eq:quartic-envelope}, every subsequence has a locally uniformly convergent subsequence. The identity theorem identifies its limit with $H$ using that imaginary interval. Hence the full cofinal sequence converges to $H$ locally uniformly on the upper half-plane. Hurwitz makes $H$ nonvanishing there, since it is not identically zero. Conjugation gives the same result below the real axis. The normalization of $H$ then gives RH.
\end{proof}

For each $M$, the measure $\nu_M=\sum_j\lambda_{M,j}b_{M,j}^2\delta_{b_{M,j}}$ converts \eqref{eq:weighted-moments} to
\[
 \int x^k\dd\nu_M(x)=Q_{k+2},\qquad0\le k\le M-2.
\]
This is a positive finite moment problem. Every available Hankel matrix
\[
 [Q_{i+j+2}]_{i,j=0}^d,\qquad [Q_{i+j+3}]_{i,j=0}^d
\]
must be positive semidefinite whenever its entries are prescribed, by integrating $p(x)^2$ and $xp(x)^2$. We do not assume a converse for a truncated system. In particular, finite feasibility through $M=13$ is not unbounded positive feasibility.

The earlier connected 305-vertex ferromagnet in \cite{R994} remains a distinct finite contribution: it matches the original standardized theta moments through degree 20 with 304 observed spins and strictly positive couplings. The present order-26 products are stronger in Taylor order but belong to a larger, not necessarily probabilistic, class. Conflating the two would misstate what either certificate proves.

\subsection{Why a finite match cannot locate all zeros}
The need for unbounded orders is not merely a technical qualification. The later note~1008 shows that a finite theta match can coexist with prescribed nonreal zeros far away, even while retaining the same simple bound on the whole plane \cite{R1008}. The construction is short enough to include here.

\begin{proposition}[Preserving a finite match while inserting four zeros]\label{prop:finite-insertion}
Suppose
\[
 P(z)=e^{-az^2}\prod_{j=1}^N(1-b_jz^2),\qquad
 a>0,\quad b_j>0,\quad a+\sum_jb_j=\tfrac12,
\]
with $\sum_jb_j^r=Q_r$ for $2\le r\le M$. Suppose $M-1$ chosen nodes are distinct, so their power-sum Jacobian is nonsingular. Fix $v>0$ and $B>0$. For every sufficiently large $T$, there is an even real entire function $G_T$ with the same Taylor coefficients through order $2M$, satisfying
\[
 |G_T(z)|\le e^{|z|^2/2}\qquad(z\in\C),
\]
whose only nonreal zeros in the physical variable, for $g_T(w)=G_T(\sqrt v\,w)$, are $\pm T\pm\ii B$.
\end{proposition}
\begin{proof}
Put
\[
 c=\frac1{v(T+\ii B)^2}=u+\ii y.
\]
As $T\to\infty$, $c\to0$. The implicit-function theorem permits a small change in the chosen real nodes such that
\[
 \sum_j\widetilde b_j^r=Q_r-2\Re(c^r),\qquad 2\le r\le M.
\]
Indeed their Jacobian has entries $r b_j^{r-1}$, a nonsingular diagonal multiple of a Vandermonde matrix. All the changed nodes remain positive. For large enough $T$, the adjusted Gaussian coefficient
\[
 \widetilde a=\tfrac12-\sum_j\widetilde b_j-2u
\]
is also positive. Define
\[
 G_T(z)=e^{-\widetilde a z^2}
   \prod_j(1-\widetilde b_jz^2)(1-cz^2)(1-\overline c z^2).
\]
The logarithmic coefficients give the stated Taylor match. The last two factors give precisely the four prescribed physical zeros; the other factors have real zeros.

For sufficiently large $T$, $u>0$ and $y^2\le u^2$. Writing $t=|z|^2$ gives
\[
 \begin{split}
 |(1-cz^2)(1-\overline c z^2)|
 &\le 1+2ut+(u^2+y^2)t^2\\
 &\le 1+2ut+2u^2t^2\le e^{2ut}.
 \end{split}
\]
The bounds for the real-node factors and the defining identity for $\widetilde a$ now give the unchanged envelope $e^{|z|^2/2}$.
\end{proof}

The source applies this argument to the nonsingular finite constructions through order $26$. It does not construct nonreal zeros of $\xi$: the functions $G_T$ are auxiliary functions. The order $M$ is fixed first, and the inserted zeros are then moved far away. This is exactly the distinction that an argument for RH must overcome.

There is even a quantitative indication of how little a finite test can see. The extra squared absolute node mass is
\[
 2|c|^2-2\Re(c^2)
 =\frac{16T^2B^2}{v^2(T^2+B^2)^4}=O_{v,B}(T^{-6}).
\]
A separate probability-law construction in \cite{R1008} retains positivity and several leading growth properties while allowing both zero geometries. Its theta application uses the finite Ising model through order $20$, not the order-$26$ entire-function construction. Neither result contradicts the all-order criterion: both show why its unbounded feasibility hypothesis cannot be replaced by finitely many successful matches.

\section{Strip-confined models: compactness and controlled negative couplings}\label{sec:strips}

The Lee--Yang real-zero class has strong closure properties; see \cite{LY52,NW19}. Merely replacing real zeros by a common horizontal strip does not preserve those properties. The counterexample and repair in \cite{R1002} are therefore part of the mathematical result, not optional cautions about numerical fitting.

\subsection{A moment-convergence counterexample}
For $a\ge0,w>0$,
\[
 f_{a,w}(z)=\frac{a+\cos(wz)}{a+1}
\]
is the characteristic function of a symmetric three-point law. If $0\le a\le1$ all its zeros are real. If $a>1$, the zeros are
\[
 \frac{(2j+1)\pi\pm\ii\arcosh(a)}{w},\qquad j\in\Z.
\]
Fix $b>0$, set $a_n=\cosh(bn)$ and $v_n=n^2/(1+a_n)$, and, for large $n$, define
\begin{equation}\label{eq:rare-jump}
 f_n(z)=\cos(\sqrt{1-v_n}\,z)\,
          \frac{a_n+\cos(nz)}{1+a_n}.
\end{equation}
These are symmetric characteristic functions of variance exactly one. All their zeros lie in $\abs{\Im z}\le b$. Every fixed moment converges to the corresponding moment of the fair sign law: the rare-jump contribution to the $2k$th moment is $n^{2k}/(1+\cosh(bn))\to0$, and $v_n\to0$. Yet, for every $y>b$,
\[
 f_n(\ii y)\ge\frac{\cosh(bn)+\cosh(ny)}{1+\cosh(bn)}\longrightarrow\infty.
\]
Thus even convergence of every fixed moment does not give local boundedness in the region where zero exclusion would be transferred.

\begin{theorem}[One exponential budget suffices]\label{thm:strip-transfer}
Let $\nu$ be the actual theta probability law with characteristic function $F_\theta(z)=\Xi(z)/\Xi(0)$. Let $\nu_m$ be symmetric probability measures with entire characteristic functions $f_m$. Suppose, for fixed $0\le b<1/2<R$ and $C<\infty$, that all zeros of $f_m$ lie in $\abs{\Im z}\le b$, that
\begin{equation}\label{eq:one-exp-budget}
 \int e^{R\abs x}\dd\nu_m(x)\le C,
\end{equation}
and that every fixed moment converges to the corresponding theta moment. Then $\Xi$ has no zeros with $\abs{\Im z}>b$. Hence ordinary zeta has no nontrivial zero with $\Re\rho>1/2+b$.
\end{theorem}
\begin{proof}
For $\abs{\Im z}<R$, $\abs{f_m(z)}\le C$, so the functions form a normal family on this connected strip. Every subsequential limit has at zero the prescribed derivatives, since local uniform convergence permits differentiation. It therefore equals $F_\theta$ near zero, and then on the strip by analytic uniqueness. This identifies every subsequential limit and proves local uniform convergence of the sequence. Hurwitz applies on each of $b<\Im z<R$ and $-R<\Im z<-b$. The limit is not identically zero because its value at zero is one. All possible nontrivial zeta zeros already correspond to $\abs{\Im z}<1/2<R$, which proves the conclusion at every height.
\end{proof}

It suffices to certify the zero exclusion on expanding horizontal windows inside $\abs{\Im z}<R$, because each compact eventually lies in such a window. If the first $L$ moments match exactly and both measures satisfy the exponential budget $C$, the elementary Taylor-tail estimate is
\[
 \sup_{\abs z\le r}\abs{f_m(z)-F_\theta(z)}
 \le2C(r/R)^{L+1},\qquad r<R.
\]
Moment mismatches add their explicit finite Taylor sum. This is a disc estimate, not a uniform assertion on an unbounded horizontal strip.

The finite three-point factors above describe a test class and the counterexample, not a successful cofinal theta construction. The separate product obstruction in \cite[strip proof]{R1001} further warns against treating products of such elementary laws as a universal approximation class. The correlated graph construction below has different hypotheses and is not excluded by that product argument.

\subsection{A finite-graph strip theorem}
Consider a finite graph without self-loops, with finite real pair couplings $J_e$ and real fields $h_i$. Under the probability law proportional to
\[
 \exp\left(\sum_{e=\{i,j\}}J_e\sigma_i\sigma_j+\sum_i h_i\sigma_i\right),
 \qquad \sigma_i\in\{-1,1\},
\]
let $Y=\sum_iw_i\sigma_i$, $w_i\ge0$, and $\varphi(z)=\E e^{\ii zY}$. Parallel edges may first be combined. The following result is the controlled-negative-coupling extension in \cite{R1003}; its proof gives the needed disk form of the classical contraction argument directly.

\begin{theorem}[Allocatable negative-edge budget]\label{thm:negative-strip}
For each negative edge $e=\{i,j\}$, choose nonnegative endpoint budgets $\alpha_{e,i},\alpha_{e,j}$ satisfying
\begin{equation}\label{eq:edge-budget}
 \tanh\alpha_{e,i}\tanh\alpha_{e,j}\ge\tanh\abs{J_e}.
\end{equation}
Suppose, at each vertex,
\begin{equation}\label{eq:vertex-budget}
 \abs{h_i}+\sum_{e\ni i,\ J_e<0}\alpha_{e,i}\le b w_i.
\end{equation}
Then every zero of the complete characteristic function $\varphi$ lies in $\abs{\Im z}\le b$. There is no restriction on the number of vertices, cycles, or positive couplings.
\end{theorem}
\begin{proof}
For an edge put $c=e^{-2J}>0$. Its two-variable polynomial, up to a nonzero positive constant, is
\[
 p_c(x,y)=1+c(x+y)+xy.
\]
If $c\le1$, it is nonzero when $\abs x,\abs y<1$. For $c<1$, the possible root $y=-(1+cx)/(c+x)$ has modulus greater than one, since
$\abs{1+cx}^2-\abs{c+x}^2=(1-c^2)(1-\abs x^2)>0$; the case $c=1$ factors.

If $c>1$, choose $r_i,r_j\in(0,1)$ with
$1+r_ir_j\ge c(r_i+r_j)$. This implies $r_i,r_j<1/c$. Minimization on a circle gives, for $\abs x\le r_i$,
\[
 \left|\frac{1+cx}{c+x}\right|
 \ge\frac{1-c\abs x}{c-\abs x}
 \ge\frac{1-cr_i}{c-r_i}\ge r_j.
\]
Thus the edge polynomial is nonzero in the open bidisk. With $r_i=e^{-2\alpha_i}$ and $r_j=e^{-2\alpha_j}$, the radius condition is equivalent to \eqref{eq:edge-budget}, because
\[
 \frac{1+r_ir_j}{r_i+r_j}
 =\frac{1+\tanh\alpha_i\tanh\alpha_j}
        {1-\tanh\alpha_i\tanh\alpha_j}.
\]

The disk contraction step is elementary. If $P(x,y)=A+Bx+Cy+Dxy$ is nonzero on $\abs x<r_1,\abs y<r_2$, then $A+Dz$ is nonzero for $\abs z<r_1r_2$. Indeed $A\ne0$, and the roots of $P(r_1t,r_2t)$ have modulus at least one. Their product gives $\abs D r_1r_2\le\abs A$; if $D=0$ the conclusion is immediate. Other variables can be held fixed throughout, so the contraction can be iterated.

Assign separate variables to edge incidences, multiply all edge polynomials, and contract the incidences at each vertex. An isolated vertex supplies $1+z_i$. The result is exactly the graph subset polynomial
\[
 \mathcal Z(\boldsymbol z)
 =\sum_{S\subseteq V}
  \left(\prod_{e:\,\abs{e\cap S}=1}e^{-2J_e}\right)\prod_{i\in S}z_i.
\]
Every physical spin configuration remains. The contraction removes only inconsistent incidence assignments introduced in the factorization. It proves nonvanishing when all $\abs{z_i}<R_i$, where
$R_i=\exp(-2\sum_{e\ni i,J_e<0}\alpha_{e,i})$. Spin reversal gives nonvanishing also when all $\abs{z_i}>1/R_i$.

Set $z_i=e^{2h_i+2\ii zw_i}$. For positive weights, \eqref{eq:vertex-budget} puts these activities strictly inside the respective disks when $\Im z>b$, and strictly outside the reciprocal disks when $\Im z<-b$. The additional exponential and normalizing factors are nonzero. A zero-weight vertex has zero field and no incident negative edge. Replacing those zero weights by positive $\eps$ and then letting $\eps\downarrow0$ gives the result by Hurwitz; the finite entire limit has value one at zero and is not identically zero.
\end{proof}

A simple symmetric allocation is
\[
 \alpha_{e,i}=\alpha_{e,j}
 =\tfrac12\arcosh(e^{2\abs{J_e}}),
\]
but asymmetric allocations are allowed. Thus a larger observable weight can pay a larger share of an incident negative-edge cost. This theorem does not assign the ordinary Lee--Yang real-zero property to an antiferromagnet; it pays an explicit nonzero strip width.

\subsection{A reference-variance bound and its continuous barrier}
Let the positive reference model have zero fields and couplings $J_e^+=\max(J_e,0)$. Put
\[
 C_{ij}=\E_+(\sigma_i\sigma_j),\qquad V_+=w^{\mathsf T}Cw,
 \qquad K_-=\sum_{J_e<0}\abs{J_e}.
\]
The correlation matrix is positive semidefinite, has diagonal one, and has nonnegative entries. The zero-field ferromagnetic Gaussian estimate is
\begin{equation}\label{eq:ferro-Gaussian}
 \E_+e^{\sum_i q_i\sigma_i}\le e^{q^{\mathsf T}Cq/2},\qquad q_i\ge0.
\end{equation}
This is the classical Lee--Yang consequence used in \cite[strip theorem]{R1001}. For example, apply the real-zero product to the finite weighted observable $\sum_iq_i\sigma_i$ and use $\log(1+x)\le x$ on its moment generating product. Arbitrary field signs can be replaced by absolute values on the right, by the nonnegative zero-field spin correlations.

\begin{proposition}[Reference variance and entire growth]\label{prop:reference-variance}
Under \eqref{eq:edge-budget}--\eqref{eq:vertex-budget}, let $v=\E Y^2$. Then
\begin{align}
 v&\ge V_+e^{-b^2V_+},\label{eq:variance-barrier}\\
 \abs{\varphi(z)}&\le
 \exp\left[V_+\left(b^2+b\abs{\Im z}+\tfrac12\abs{\Im z}^2\right)\right].
 \label{eq:graph-envelope}
\end{align}
Here $v$ is a second moment; no centering is needed for these inequalities.
\end{proposition}
\begin{proof}
Write $d_i=\sum_{e\ni i,J_e<0}\alpha_{e,i}$ and $q_i=\abs{h_i}$. The edge condition implies
\[
 \abs{J_e}\le\tfrac12\{
 \log\cosh(\alpha_{e,i}+\alpha_{e,j})
 -\log\cosh(\alpha_{e,i}-\alpha_{e,j})\}
 \le\alpha_{e,i}\alpha_{e,j},
\]
where the last bound integrates $\tanh t\le t$. Hence
\[
 2K_-\le\sum_i d_i^2\le d^{\mathsf T}Cd,
 \qquad 2K_-+\tfrac12q^{\mathsf T}Cq\le b^2V_+.
\]
For the second inequality use $d+q\le bw$ componentwise and nonnegative entries of $C$.

Let $A=\sum_{J_e<0}J_e\sigma_i\sigma_j$ and $H=\sum_i h_i\sigma_i$. Global spin reversal of the reference law gives
$\E_+[Y^2e^{A+H}]=\E_+[Y^2e^A\cosh H]\ge e^{-K_-}V_+$. Its normalizer is at most $e^{K_-+q^{\mathsf T}Cq/2}$ by \eqref{eq:ferro-Gaussian}. Division proves \eqref{eq:variance-barrier}.

The normalizer is also at least $e^{-K_-}\E_+e^H\ge e^{-K_-}$. With $y=\abs{\Im z}$, the absolute numerator is at most
$e^{K_-}\exp((q+yw)^{\mathsf T}C(q+yw)/2)$. Use the previous budget and $q\le bw$ to obtain \eqref{eq:graph-envelope}.
\end{proof}

A continuous deformation preserving $v$, the band budget $b$, and continuity of $V_+$ cannot escape its small-reference-variance component when it begins at a zero-field ferromagnet. If $b^2v<1/e$, the inequality $v\ge V_+e^{-b^2V_+}$ excludes the open interval between the two roots of $Ve^{-b^2V}=v$. Continuity prevents crossing that interval.

For the physical theta normalization, take $b=1/3$ and use $v_\theta<4621/100000$. The explicit barrier $r=2323/50000=0.04646$ satisfies
\[
 r-r^2/9-4621/100000=\frac{228671}{22500000000}>0.
\]
Since $e^{-x}\ge1-x$, a first crossing $V_+=r$ would contradict \eqref{eq:variance-barrier}. Thus $V_+<r$ throughout such a continuous component. Graph changes are allowed when they preserve this continuity, for example the insertion of zero edges or zero-weight vertices.

A cofinal sequence of \emph{symmetric}, theta-matching models in this component would have the common entire envelope \eqref{eq:graph-envelope}. Montel, coefficient identification, and Hurwitz would then give the ordinary $5/6$ boundary. No such cofinal sequence, or negative-coupling theta-matching finite model, is supplied by \cite{R1003}. The contribution is an enlarged admissible class with a paid strip and compactness budget, not an achieved zero-free region.

\section{Scale filters and a determinant that retains the primes}\label{sec:detectors}

The two constructions in this section ask a related question: can we change the arithmetic expression without erasing the zero we hope to rule out?

A scale difference can remove a known main term, but may also remove certain oscillations. Using two different scales avoids a common nonreal blind spot. A determinant provides a different representation: its logarithm is an exact weighted sum of prime powers. In both cases the representation is useful because its Mellin transform still detects every relevant zero. The work in \cite{R993,R1000,R1002} establishes these identities and bounds their approximation errors. An upper bound strong enough to exclude new zeros remains a separate task.

\subsection{Two scale differences without a common nonreal blind spot}
For real $\lambda$ and $h>0$, define
\[
 \Delta_{\lambda,h}F(X)=F(X)-e^{-\lambda h}F(e^hX).
\]
Its $m$th power is the complete finite combination
\begin{equation}\label{eq:scale-difference}
 \Delta_{\lambda,h}^mF(X)=\sum_{j=0}^m(-1)^j\binom mj
                                   e^{-\lambda hj}F(e^{jh}X).
\end{equation}
It annihilates $X^\lambda P(\log X)$ whenever $\deg P<m$, by the ordinary finite-difference identity. For a convergent source $F(X)=\sum_n a_nW(n/X)$, the operation keeps every coefficient and changes the profile to
$\sum_{j=0}^m(-1)^j\binom mj e^{-\lambda hj}W(e^{-jh}x)$. Ideal norms can replace $n$ without changing this identity.

With Mellin convention $\mathcal MF(s)=\int_0^\infty F(X)X^{-s}\dd X/X$, substitution gives, on an initial convergence strip,
\begin{equation}\label{eq:notch-multiplier}
 \mathcal M(\Delta_{\lambda,h}^mF)(s)
 =[1-e^{h(s-\lambda)}]^m\mathcal MF(s).
\end{equation}
One multiplier vanishes on the vertical lattice $\lambda+2\pi\ii\Z/h$. Taking $h=1$ and $h=\sqrt2$ leaves only $s=\lambda$ in their common zero set.

\begin{proposition}[Joint preservation and polynomial height cost]\label{prop:paired-notches}
Suppose a fixed-source Mellin transform is $H(s)/\zeta(s)$, or $H(s)[-\zeta'(s)/\zeta(s)]$, up to a holomorphic term, and $H$ is nonzero at every target zero. Fix $\lambda\in(0,1)$ and $m\ge1$. If both transforms in \eqref{eq:notch-multiplier}, at $h=1,\sqrt2$, are holomorphic on $\Re s>\theta$, then there is no nontrivial zeta zero there. On $\abs{\Re s-\lambda}\le A$, $\abs{\Im s}\ge2\pi$,
\begin{equation}\label{eq:paired-lower}
 \max_{h\in\{1,\sqrt2\}}\abs{1-e^{h(s-\lambda)}}
 \ge\frac{e^{-\sqrt2 A/2}}{3(1+\abs{\Im s})}.
\end{equation}
\end{proposition}
\begin{proof}
At any nonreal zero, at least one multiplier is nonzero and preserves the reciprocal pole order, or the logarithmic-derivative residue with its multiplicity. The common real point is not a zeta zero: for $0<s<1$, the alternating eta series is positive while $1-2^{1-s}<0$.

For the quantitative bound, put $u=\Im s/(2\pi)$, and choose nearest integers $n$ to $u$ and $k$ to $\sqrt2u$. With $e=\max(\norm u_{\R/\Z},\norm{\sqrt2u}_{\R/\Z})$ and $\abs u\ge1$, $n\ne0$ and
\[
 1\le\abs{2n^2-k^2}
 \le(\sqrt2+1)e\{2\sqrt2\abs u+(\sqrt2+1)/2\}
 \le12e(1+\abs u).
\]
Combine this with $\abs{1-e^{2\pi\ii v}}\ge4\norm v_{\R/\Z}$ and
\[
 \abs{1-e^{a+\ii b}}^2=(e^a-1)^2+e^a\abs{1-e^{\ii b}}^2.
\]
The real-part range costs at most $e^{-\sqrt2 A/2}$, proving \eqref{eq:paired-lower}. For the $m$th differences raise the bound to the power $m$.
\end{proof}

The same real-point argument applies to the \emph{principal} Dedekind zeta functions of $\Q(\sqrt{-3})$ and $\Q(\ii)$. Indeed, for every real $s>0$,
\[
 L(s,\chi_{-3})=\sum_{n\ge0}[(3n+1)^{-s}-(3n+2)^{-s}]>0,
\]
\[
 L(s,\chi_{-4})=\sum_{n\ge0}[(4n+1)^{-s}-(4n+3)^{-s}]>0.
\]
The paired series converge normally on $\Re s>0$ by the mean-value estimate on each difference, and agree with the Dirichlet series initially. Equation~\eqref{eq:quadratic-factorization} proves the required real nonvanishing. This does not supply such a sign argument for every quadratic field or every auxiliary character.

Applying \eqref{eq:scale-difference} to a full physical theta expression still requires checking its linked rescalings and every remainder. The old triangle bounds reproduce the old exponent. In particular, removing a real term in a first Gauss-sum mean does not remove the whole exceptional-frequency covariance in its squared norm.

\subsection{A rational positive matrix for an exact prime-power integer}
Let $S_n$ be the cyclic permutation matrix, $L_n=2I-S_n-S_n^{\mathsf T}$, and let
\[
 (P_n)_{ij}=\frac{c_n(i-j)}n,\qquad 0\le i,j<n.
\]
Finite Fourier inversion identifies $P_n$ with the rational orthogonal projector onto all primitive $n$th Fourier modes. It commutes with $L_n$, and its rank is $\phi(n)$. Define
\begin{equation}\label{eq:positive-matrix}
 H_n(t)=I-P_n+P_n(L_n+tI),\qquad
 H_N(t)=\bigoplus_{n=2}^N H_n(t)^{\oplus(N-n)},
 \qquad V_N(t)=\tfrac12\log\det H_N(t).
\end{equation}
An empty direct sum has determinant one. Complementary modes have eigenvalue one, while the primitive modes have eigenvalues
$t+4\sin^2(\pi a/n)$, $(a,n)=1$. All these matrices are positive definite for $t\ge0$, and have rational entries for rational $t$.

\begin{proposition}\label{prop:exact-det}
For $n\ge2$,
\[
 D_n(t):=\det H_n(t)
 =\operatorname{Res}_x(\Phi_n(x),x^2-(2+t)x+1)\in\Z[t],
\]
and
\begin{equation}\label{eq:prime-det}
 D_n(0)=\Phi_n(1)^2=e^{2\Lambda(n)},\qquad
 \det H_N(0)=A_N^2,
 \quad A_N=\prod_{p^j\le N}p^{N-p^j}.
\end{equation}
\end{proposition}
\begin{proof}
Evaluate the quadratic at each primitive root of unity. The product of the factors $-\zeta$ is one, leaving exactly the product of the primitive eigenvalues. The resultant is integral and its polynomial factors in $t$ all have positive constant terms. The classical identity
$\Phi_n(1)=\prod_{d\mid n}d^{\mu(n/d)}=e^{\Lambda(n)}$
includes every prime power; see, for example, the cyclotomic discussion in \cite{AK22}. Multiplying with the complete multiplicities $N-n$ gives \eqref{eq:prime-det}.
\end{proof}

Write
\[
 C(x)=\sum_{n\le x}(x-n)\Lambda(n),\qquad C(N)=\log A_N.
\]
This is also $\log\prod_{r<N}\operatorname{lcm}(1,\ldots,r)$ at integer $N$. The identity is a classical Riesz-mean quantity in a matrix representation; it is not a spectrum of Riemann zeros.

\subsection{Regularizing the determinant and controlling the error}
Continuously at zero, put
\[
 f(u)=\log\frac{\sinh(u/2)}{u/2},\quad
 F(v)=\int_0^1(1-y)f(vy)\dd y,\quad
 \mathcal B(b)=\sum_{m\ge1}\frac{\mu(m)}mF(b/m).
\]
Euler's product for $\sinh$ gives, for every $u\ge0$,
\[
 0\le f(u)\le u^2/24,\quad 0\le f'(u)\le u/12,\quad f''(u)\ge0.
\]
Thus the series is absolutely convergent, with complete tail at $m>K$ bounded by $b^2/(576K^2)$. Rearranging the positive logarithmic product gives
\begin{equation}\label{eq:B-positive}
 \mathcal B(b)=\sum_{a\ge1}\frac{\phi(a)}a
 \int_0^1(1-y)\log\left(1+\frac{b^2y^2}{4\pi^2a^2}\right)\dd y.
\end{equation}
The identity $\sum_{m\mid a}\mu(m)/m=\phi(a)/a$ supplies the coefficient. Absolute convergence under the quadratic majorant justifies the rearrangement.

\begin{theorem}[Every-mass comparison]\label{thm:all-mass}
For every integer $N\ge2$ and every real $b\ge0$,
\begin{equation}\label{eq:all-mass}
 \left|V_N\bigl(2(\cosh(b/N)-1)\bigr)-C(N)-N^2\mathcal B(b)\right|
 \le b^2\left(\frac{h_N}{72}+\frac1{24N}+\frac1{576}\right).
\end{equation}
Equivalently, at any mass $t\ge0$ the exact counterterm uses $b=N\arcosh(1+t/2)$.
\end{theorem}
\begin{proof}
With $h=b/N$, the complete cyclotomic divisor identity and
$\log(1-e^{-u})=\log u-u/2+f(u)$ give
\[
 V_N\bigl(2(\cosh h-1)\bigr)-C(N)
 =\sum_{n=2}^N(N-n)\sum_{d\mid n}\mu(n/d)f(hd).
\]
Set $g_m(x)=(N-mx)f(bx/N)$ on $[0,N/m]$ and extend by zero. Reindexing the preceding expression gives
\[
 \sum_{m\le N}\mu(m)\sum_{d\in\Z}g_m(d)-(N-1)f(b/N).
\]
The last term removes the unit $n=1$, and
$\int_\R g_m(x)\dd x=N^2F(b/m)/m$.

The endpoint derivative jump at $N/m$ must be included. The total variation of the second distributional derivative is bounded by
\[
 \operatorname{TV}(\mathrm dg_m')\le4m f(b/m)\le\frac{b^2}{6m}.
\]
Indeed the two interior absolute contributions integrate to $2mf(b/m)$ and $mf(b/m)$, while the endpoint jump contributes a further $mf(b/m)$. Twice integrating the periodic Bernoulli polynomial, whose bound is $\abs{B_2(\{x\})}\le1/6$, gives
\[
 \left|\sum_{d\in\Z}g_m(d)-\int_\R g_m(x)\dd x\right|
 \le\frac1{12}\operatorname{TV}(\mathrm dg_m')\le\frac{b^2}{72m}.
\]
Sum with $\abs{\mu(m)}\le1$, pay the full $m>N$ tail by $b^2/576$, and use $(N-1)f(b/N)\le b^2/(24N)$. This proves \eqref{eq:all-mass}. No step requires $b\le1$.
\end{proof}

The earlier small rational-mass statement in \cite{R1000} follows, for $0\le c\le1$, by setting $b=2N\arsinh(c/(2N))$ and comparing $b$ with $c$:
\[
 \abs{V_N(c^2/N^2)-C(N)-N^2\mathcal B(c)}
 \le\frac{c^2(1+\log N)}{24}.
\]
For a growing regularizer, retaining the exact $b$ avoids an additional, inadequately small-parameter correction.

Fix $1/2\le\theta<1$ and choose
\[
 c_N=\lceil N^{(1+\theta)/2}\rceil,\quad
 t_N=c_N^2/N^2,\quad b_N=2N\arsinh(c_N/(2N)).
\]
Then $N^{\theta-1}\le t_N\le1$. For $N\ge3$, all eigenvalues of $H_N(t_N)$ lie in $[t_N,5]$, so
\begin{equation}\label{eq:conditioning}
 \cond(H_N(t_N))\le5N^{1-\theta},\qquad
 V_N(t_N)=C(N)+N^2\mathcal B(b_N)+O(N^{1+\theta}\log(2N)).
\end{equation}
At the RH-calibrated scale the condition-number bound is $O(N^{1/2})$, and at the scale calibrated to $5/6$ it is $O(N^{1/6})$. Neither exponent is an achieved zeta boundary.

The larger regularization also creates a larger bulk. Since $f(u)/u\to1/2$ and $0\le f(u)\le u/2$, dominated convergence gives
\[
 \frac{\mathcal B(b)}b\longrightarrow\frac1{12\zeta(2)}.
\]
Thus $N^2\mathcal B(b_N)$ has order $N^{(5+\theta)/2}$, larger than the prime main term. Its leading linear asymptotic has not been shown sufficiently accurate for a zero-free deduction; the exact counterterm in \eqref{eq:conditioning} remains necessary.

If $D_N$ is the matrix dimension, $D_N\le N^3/6$, and the eigenvalue form of Frullani's identity yields
\[
 V_N(t)=\frac12\int_0^\infty
       \{D_Ne^{-u}-\Tr(e^{-uH_N(t)})\}\frac{\dd u}{u}.
\]
For $0<t\le1$, the absolute tail after $T$ is at most $D_Ne^{-tT}/(tT)$. Taking $T=K\log N/t_N$ pays the entire longer-time tail by
$N^{3-K}/(6K\log N)$ and uses only $T\le K N^{1-\theta}\log N$. The centered short-time trace estimate remains unproved.

\subsection{The missing bound on the determinant}
Let $E(x)=C(x)-(x-1)^2/2$. Direct integration on $\Re s>1$ gives
\begin{equation}\label{eq:det-Mellin}
 \int_1^\infty E(x)x^{-s-2}\dd x
 =\frac{-\zeta'(s)/\zeta(s)-1/(s-1)}{s(s+1)}.
\end{equation}
At each nontrivial zero $\rho$, the right side has residue $-m_\rho/[\rho(\rho+1)]\ne0$. It is holomorphic at every real $s>0$; the apparent pole at one has been removed.

If the additional one-sided estimate
\begin{equation}\label{eq:open-height}
 V_N(t_N)\le\frac{(N-1)^2}{2}+N^2\mathcal B(b_N)
                          +O_\eps(N^{1+\theta+\eps})
\end{equation}
were proved for every $\eps>0$ and all integers $N$, then \eqref{eq:conditioning} would give $E(N)\ll_\eps N^{1+\theta+\eps}$. Linear interpolation of $C$ shows that $E(x)$ exceeds its endpoint interpolation by at most $1/8$, so the estimate extends to real $x$.

For completeness, the one-sided implication is not an unjustified absolute-value bound. Beyond a fixed initial interval, $Kx^{1+\theta+\eps}-E(x)$ is nonnegative. Landau's positivity lemma says that a finite convergence abscissa of its Laplace transform is a singularity at the real abscissa. The meromorphic expression from \eqref{eq:det-Mellin} is analytic at every real $s>\theta+\eps$, so the convergence abscissa is at most $\theta+\eps$. Subtracting the explicit $K$ term then makes \eqref{eq:det-Mellin} holomorphic throughout that half-plane. Its nonzero zero residues exclude every zero there. Choosing $\eps$ separately at each point proves exclusion to the right of $\theta$.

The positivity lemma follows directly by expanding about a point just to the right of an allegedly regular real abscissa: the Taylor coefficients of the translated Laplace transform are nonnegative weighted moments, so monotone convergence would extend convergence to its left, a contradiction. No rightmost-zero or simplicity assumption is needed. At $\theta=1/2$, the classical RH bound for $\psi(x)-x$, followed by integration, supplies the converse; the missing estimate is therefore RH-strength. Positive definiteness of \eqref{eq:positive-matrix} does not prove \eqref{eq:open-height}.

\section{Source-qualified boundary proposals and an exact failed comparison}\label{sec:boundaries}

The components above must be distinguished from a zero-free theorem. This section records the strongest relevant broad-field proposal, the separate quadratic prime-bias proposal, and the failed ordinary-$0.87$ comparison. Their scopes are not interchangeable. In particular, a proposed boundary for every number field need not improve an already reported, narrower ordinary-zeta boundary.

\subsection{The proposed all-number-field comparison}
The frozen composition \cite{R977} proposes nonvanishing for every Dedekind zeta function on
\begin{equation}\label{eq:271}
 \Re s>271/300,\qquad s\ne1,
\end{equation}
and the same strict region for finite-order Hecke $L$-functions. It explicitly depends on new analytic adapters in preceding research branches. This manuscript does not independently reconstruct those adapters, and therefore records the conclusion as source-qualified rather than as a newly certified theorem.

The main physical input is the proposed short-support reflection estimate
\begin{equation}\label{eq:short-reflection-contract}
 \abs{I(X,Y,Z)}\ll_\eta Z^\eta
       \left[(XY+(XY)^2/Z)(X^{-1}+Y^{-1})\right]^{1/2}.
\end{equation}
It retains the diagonal $XY$, the full cube variable, all local auxiliary ideals, and effective lengths below one. Its use is not justified by substituting new parameters into a long-support corollary. The other main adapter is an enlarged holomorphy domain for the exact correction factor in the Poisson expansion
\[
 X^{1/2-z}Z^{\xi+z-1}Y^{w-1}q_u^{-z}\zeta_F^S(6z)
          \frac{L_F^S(w,\chi_u)H_u(\xi,w,z)}{L_F^S(\xi,\overline\chi_u)}.
\]
For $x=\Re\xi$, $y=\Re w$, the adopted domain is
\begin{gather}\label{eq:correction-domain}
 x,y,z>0,\quad x\ge1/2,\quad 6z+y>1,\quad x+6z>1,\quad x+y>1,\notag\\
 x+y+6z>2,\qquad 6x+6z>5.
\end{gather}
The finite-$S$, ramified, and full-norm character conventions are essential. The source also adopts a complete second numerator moment and selected-character density estimates with polynomial height dependence. None is an estimate for an unrestricted reciprocal moment.

The new use of these inputs is to reflect a numerator into the \emph{same} selected zero-free rectangle already used for its denominator. At a selected threshold $\tau$, the contour is $x=\tau+3e$, $y=1-\tau-e$, $z=1/3$. The functional equation changes the numerator to $L(\tau+e-\ii t,\overline\chi_u)$. Both signs of its analytic logarithm are then controlled in the buffered rectangle. The numerator ordinates and denominator ordinates are separately truncated; the claim does not extend this bound to arbitrary height tails.

For the exact exponent comparison, take
\[
 a=21/50,\qquad h=1-a=29/50,\qquad \vartheta=a/2=21/100,
 \qquad X=Y=Z^a.
\]
There are 851 thresholds $\sigma_j=63/100+j/2500$, $0\le j\le850$. A middle class has density at $\sigma_j$ and zero-freeness at $\tau=\sigma_{j+1}$. Its central exponent, before the reserved small losses, is
\begin{equation}\label{eq:middle-exponent}
 c_j=h(2\tau+g_j)+a-3/2.
\end{equation}
On $7/10\le\sigma\le9/10$, the adopted selected-fourth density is
\[
 g_4(\sigma)=\frac{468\sigma^2-741\sigma+271}
                  {6(36\sigma^2-75\sigma+31)}.
\]
Elsewhere the source uses $g_{\mathrm{DF}}=1-\delta_{\mathrm{DF}}$, with $q=2\sigma-1$ and
\[
 \delta_{\mathrm{DF}}=\begin{cases}
 q/(4-6q),&\sigma\le2/3,\\
 \dfrac{q\{36q-6(6\sigma-2)+3\}}
 {36q(3-2\sigma)-24\sigma+2},&2/3\le\sigma\le11/12,\\[3pt]
 \dfrac{8\sigma^2-10\sigma+3}{1-24(1-\sigma)^2},&\sigma\ge11/12.
 \end{cases}
\]
These are the exact analytic contracts and rational formulas in \cite[\texttt{DENSITY.md}]{R977}, not independently improved density results here. Their 850 middle-class instances give
\begin{equation}\label{eq:271-margin}
 \min_{0\le j<850}(\vartheta-c_j)
     =\frac{96266833}{858608062500}>0.
\end{equation}
The minimum occurs at $j=411$, $\sigma_j=993/1250$. The manuscript's exact-arithmetic check independently reproduces these finite substitutions; it does not prove the analytic density bounds.

The source fixes all selected sets and the common height exponent before allocating contour losses, separately treats the good and terminal classes, retains horizontal connectors, and sums the infinite frequency tail on its original absolute-convergence lines. Its principal detector is
\[
 f_a(Z)=\frac1{2\pi\ii}\int_{(2)}
 Z^{s-b_a}e^{(s-5/6)^2}\frac{H(s)}{\zeta_F^S(s)}\dd s,
 \qquad b_a=5/6-a/3,
\]
with inherited $H$ holomorphic and nonzero on $\Re s>2/3$. The complete proposed comparison is
\[
 I(Z^a,Z^a,Z)=c_0f_a(Z)+O_e(Z^{\vartheta+8e}),\qquad c_0\ne0.
\]
Together with \eqref{eq:short-reflection-contract}, it gives the stated detector growth and hence the boundary
\[
 b_a+\vartheta=\frac56+\frac a6=\frac{271}{300}.
\]
This last Mellin deduction is straightforward once the full comparison is valid. The substantive review obligations lie in the short-support reflection, correction domain, selected-family adapters, and complete contour accounting. Passing from $F=K(\mu_{12})$ to $K$, and then to finite-order Hecke factors, uses the corresponding abelian factorizations. Infinite-order angular characters are not included.

The lower-input variant in the same source gives $19/21$ using the original density and second numerator moment. These broad-field proposals are not a new ordinary-zeta improvement in this manuscript. The project's external benchmark record \cite{R1002} reports Liu's value approximately $0.874957019420098946$, which is numerically below $271/300$. That reported external theorem is not an analytic premise of Sections~\ref{sec:quadratic}--\ref{sec:detectors}.

\subsection{Independent-jet prime bias: a separate proposed consequence}
The determinant packet \cite{R993} gives a different proposed arithmetic statement. For a primitive nonprincipal real character $\chi$ of conductor $q$, it states
\begin{equation}\label{eq:inert-bias}
 \sum_{p\le R,\ \chi(p)=-1}\frac{\log p}{p}
 \le\frac12\log R+\frac14\log q+O(1+\log\log(3R)),
 \qquad R\ge\max(q,R_0),
\end{equation}
with absolute constants. The derivative depth and grid side length are independent parameters, rather than being identified as in the imported note. Its rank condition for $n$ points with an incidence bound $M\deg F$ is
\[
 Jn>M(R+S)+2RS+(J-M)_+S.
\]
A discriminant-balanced quadratic rectangle and $J=S=\lfloor R/\lceil\log(2R)\rceil\rfloor$ then balance the integral Gram height against inert-prime divisibility. Ordinary derivatives are essential: the $p!$ factors preserve divisibility even when the derivative order reaches $p$. Hasse derivatives cannot be substituted.

The companion in \cite{R1000} extends the stated range to
\[
 R\ge64\sqrt q\,\lceil\log(2R)\rceil^2
\]
with an explicit error $4\log\lceil\log(2R)\rceil+30$, retaining the discriminant version of the leading terms. The source gives a positive Hadamard zero budget at conductor-scaled low heights. These are not fixed all-height complex zero-free half-planes.

The real-zero implication of \eqref{eq:inert-bias} illustrates its distinct scope. Put $\ell=\log q$. Its Laplace transfer, retaining small primes, ramified primes, and all higher prime powers, gives
\[
 -\frac{L'}L(1+t,\chi)
 \ge-\ell/2-O\bigl(t\ell^2+\log(2+1/t)+\log(2+\ell)\bigr).
\]
The positive Hadamard zero sum gives, at a real zero $\beta$,
\[
 -\frac{L'}L(1+t,\chi)\le\ell/2-\frac1{t+1-\beta}+O(1).
\]
Taking $t=\ell^{-2}$ yields the proposed asymptotic consequence
\[
 (1-\beta)\log q\ge1-O\left(\frac{\log\log q}{\log q}\right).
\]
This implication is recorded subject to the determinant and transfer proof in \cite{R993}; no numerical conductor threshold or independently reviewed absolute constant is supplied here. It is not used in the preceding sections. The circulated motivating note is attributed in the source to Ben Green but its provenance remains unauthenticated; the printed identifier is \texttt{hexagon:2610.00311v1}, not an arXiv identifier. We do not list it as an authenticated public paper or reproduce its PDF.

\subsection{An exact obstruction inside the attempted 0.87 estimate system}
The failed comparison in \cite{R1004} is useful precisely because it identifies an error in sufficient estimates without asserting an error in the target conjecture. Adopt the balanced comparison and witness-capacity formulas of the earlier audit \cite{R934}. Set
\[
 B=87/100,\quad \ell_0=14/75,\quad\delta=7/18,\quad x=1/2,
 \quad\alpha=5/6,\quad a=(1+\delta)/2=25/36.
\]
The auxiliary zero class $a$ is not the location of an asserted actual zero. The relaxed direct requirement $11/12-\ell/4\le B$ forces $\ell\ge\ell_0$. Allow the larger triangle $b\ge0$, $b+\ell\le1$, and set
$\ell_y=(1+b-\ell)/2$, $h=(1+b+3\ell)/2$.

For a moment parameter $\kappa$, define
\begin{gather*}
 c=2/(3\kappa),\quad D=3-(1+c)x,\quad P=(2-cx)(1-x),\quad
 J=(\alpha-\delta)D+\delta P,\\
 t=1+\delta P/(2J),\quad
 r=\{(2-cx)t-x(1-c/2)\}/D,\quad m=t-r,\\
 z_M=(1-r)/2,\quad z_S=(1-2m)/(6\kappa),\quad
 R=1-\delta+\frac{(\alpha-\delta)\delta P}{2J}.
\end{gather*}
The raw endpoint exponent relative to the fixed target is
\begin{equation}\label{eq:failed-E}
 E=a-a\ell_y-\ell/2+x\delta\ell
                     +h(R+\delta/2-1/6)-B.
\end{equation}
At $\kappa=3/4$, exact arithmetic gives
\begin{gather*}
 t=443/394,\quad r=141/197,\quad m=161/394,\\
 z_M=28/197,\quad z_S=8/197,\quad R=2363/3546.
\end{gather*}
At $\ell=\ell_0$,
\[
 E=521/19700-b/7092\ge7003/265950>1/40.
\]
This is a certificate on the entire enlarged triangle, not merely its bottom edge. Put $v=\ell-\ell_0\ge0$ and $u=1-\ell-b\ge0$. Then
\[
 E=E_{\min}+(E_\ell-E_b)v-E_bu,
 \qquad E_\ell=640/591,\quad E_b=-1/7092,
\]
and both added coefficients are nonnegative.

Even granting the more favorable parameter $\kappa=37/50$ gives
\begin{equation}\label{eq:failed-positive-gap}
 R=2075/3114,\quad E(\ell_0,b)=457/17300-b/6228,
 \quad E_{\min}=6139/233550>1/40.
\end{equation}
Within this formula $R$ increases with $\kappa$. At $b=1/8$, the same class would need $R\le214/337$, leaving the exact count-exponent shortfall
\[
 \frac{2075}{3114}-\frac{214}{337}
 =\frac{32879}{1049418}.
\]
The raw sextic-kernel diagnostic is also exact: for
$k(v)=\max(1,v,5/6+v/3,1/3+5v/6)$,
\[
 \min_{v\ge0}\{k(v)-(7/18)v\}=7/9,
\]
attained at $v=1$. This isolated calculation is not an impossibility theorem for every application of the sieve.

Equations~\eqref{eq:failed-E}--\eqref{eq:failed-positive-gap} concern upper-estimate budgets. Extra side constraints can only restrict the parameter triangle, and logarithmic improvements cannot absorb its positive power gap. A different signed physical estimate could change the comparison. The computation supplies no lower bound on the actual arithmetic residual and no evidence for an off-critical zero.

\section{Verification, remaining estimates, and interpretation}\label{sec:verification}

\subsection{Three different forms of evidence}
A written proof, an accepting finite computation, and a byte-identical publication answer different questions. A proof must justify its infinite sums, analytic continuations, and uniformity. A finite checker can establish algebraic identities, a directed finite inequality, or an existence certificate on an explicitly covered parameter cube. A manifest or Git object identity establishes which bytes were tested or published. None substitutes for either of the others.

The theta certificates in Section~\ref{sec:theta} are genuine finite existence claims in the source: their accepting argument encloses the full analytic moment source and the full finite parameter cube. By contrast, many arithmetic checkers test only finite algebra attached to an infinite proof. The numbers of predicates in their reports include repeated domain and setup guards. They are not numbers of independent theorems, prime ranges, or zeta zeros verified.

During preparation of this manuscript, the original local packets for \cite{R1000,R1002,R1004} were replayed without editing their scientific files. The complete positive-determinant replay reproduced its ordinary and optimized 3,370-predicate reports and all eight designated corrupted executions. The five-route replay authenticated fourteen manifested payloads, reproduced 20,279 predicates per mode, and rejected twelve designated corruptions. Both completed with empty standard error. The failed-$0.87$ audit reproduced its archived report byte for byte in both isolated ordinary and optimized Python. Its report is an exponent-model calculation, not a zeta-zero computation.

A separate small manuscript checker recomputes the 850 rational middle-class comparisons in \eqref{eq:271-margin}, the scalar constants of the failed comparison, the displayed residual-length ranges, the conditional recurrence exponents, the native frequency cutoff, the graph barrier, and the inequalities using the rounded bounds in Table~\ref{tab:theta-certificates}. Rechecking those rounded bounds does not reconstruct the theta source or prove its enclosure. The full \cite{R1003} certificate campaigns, the \cite{R989} analytic contracts, and the upstream reflection and local-domain arguments in \cite{R977} were \emph{not} newly executed or independently certified by compiling this manuscript. Their exact pins and reading scopes are part of the provenance ledger.

Some source packets record separate AI-agent reviews within their originating session. Those reports remain source records; they are not external human referee reports, and they are not rebranded here as new independent reviews. This manuscript was synthesized with substantial AI assistance, including source reading, mathematical exposition, code, and typesetting. No independent mathematical referee, Lean proof, repository-wide validation, actual $L$-value census, or zeta-zero computation is claimed for its preparation. Human author approval, mathematical review, and arXiv submission remain separate actions.

The distributed LaTeX source uses a conventional \texttt{amsart} build, fixed dates, local section files, and a pre-generated bibliography. It requires no shell escape or online data access to compile. The proof-source ledger and build instructions accompany the repository version. The arXiv source archive contains the necessary source files rather than the resulting PDF or transient build products. A manuscript compile and visual inspection verify presentation, not mathematical correctness.

\subsection{The estimates not supplied}
\paragraph{The ordinary scalar residual.}
Equation~\eqref{eq:remaining-scalar} removes the auxiliary field variance without removing $\abs{\zeta(s)M_Z(s)-1}^2$. A strict upper bound below one at every point in a desired region would exclude ordinary zeros, but the source family approximation does not provide it. Similarly, the signed field mean in Theorem~\ref{thm:field-average} has not been bounded at the power required by its zero-free consumer.

\paragraph{The native arithmetic contraction.}
Equations~\eqref{eq:rough-comparison} and \eqref{eq:frequency-tail-general} pay specific complete sectors. The remaining large-core or physically restricted low-frequency amplitude still needs an arithmetic estimate with a genuinely subquadratic power of the old energy, or another estimate of comparable strength. The synthetic counterexample shows why generic coefficient bounds and balance identities alone cannot supply the proposed contraction.

\paragraph{Unbounded theta feasibility.}
The finite certificates establish matching through order 26. Theorem~\ref{thm:weighted-canonical} would use positive weighted solutions at an unbounded set of orders. Neither two successful root-extension steps nor finitely many positive Hankel matrices establishes that premise. In the strip programme, controlled fields and negative edges enlarge the admissible class but have not produced a cofinal theta-matching sequence.

\paragraph{A signed filtered or determinant estimate.}
The two incommensurate filters preserve every nonreal target pole, but no improved full physical remainder estimate is proved. The all-mass determinant theorem reduces conditioning and pays the complete long heat tail, but not the centered short-time trace or height. These are different representations of a remaining arithmetic task, not automatic deductions from source identification.

\paragraph{Independent review of the broad-field proposal.}
The displayed $271/300$ composition retains the scope and dependencies of \cite{R977}. Its review requires reconstructing the short-support estimate, ramified local correction, selected-character interfaces, and all contour tails. Finite exponent margins, including the freshly recomputed \eqref{eq:271-margin}, cannot validate those analytic premises.

\subsection{Conclusion}
The recent contributions are not exhausted by a sequence of decimals near one. They include a uniform actual-field coefficient comparison; a full-$L$ approximation and an exact auxiliary-variance optimization; a two-polynomial transfer for capped reciprocal laws; complete denominator, core, and frequency-sector bounds; finite real-zero approximants through order 26; a less restrictive weighted canonical criterion; a controlled negative-coupling strip class; and explicitly regularized positive determinants. The failed-$0.87$ comparison is also part of the result record because it identifies a precise missing power saving in a specified estimate system.

What these contributions share is a requirement of source fidelity: common factors, square aliases, ramified zeros, prime powers, complete cross terms, physical cutoffs, and limit quantifiers are retained rather than discarded when inconvenient. They provide component theorems and more explicit targets. They do not yet provide the all-scale signed estimate or unbounded positive construction needed for a new global ordinary-zeta zero-free region or RH.

\appendix
\section{Sources and related results}\label{app:provenance}

This paper is based on frozen versions of the research notes, not on an assumed single merged proof. Table~\ref{tab:pins} gives a short source map. The bibliography links to the immutable versions; the accompanying \texttt{SOURCE\_LOCK.json} gives the full commit identifiers, principal paths, recorded file hashes, and the scope of the reading used for this revision. The longer technical proofs and finite parameter files remain available there.

\small
\begin{longtable}{@{}p{1.5cm}p{2.4cm}>{\raggedright\arraybackslash}p{10.7cm}@{}}
\caption{Research notes used in the paper. A version identifier records the source, not a verdict on its correctness.}\label{tab:pins}\\
\toprule Note & Version & Principal use\\
\midrule\endfirsthead
\toprule Note & Version & Principal use\\\midrule\endhead
\bottomrule\endfoot
934 & \texttt{3092feee25c5} & Earlier exponent model and its limits.\\[3pt]
957 & \texttt{ec54cb6b0029} & Stronger density; exact family matching; 23/24 proposal\\[3pt]
958 & \texttt{e6dfa07d07ed} & Finite zero strata and the 25/27 proposal\\[3pt]
959 & \texttt{7b845294207d} & Second numerator mean and the 11/12 proposal\\[3pt]
965 & \texttt{d58e3282a937} & Two-term short-support estimate; 91/100 proposal\\[3pt]
966 & \texttt{8d94494698ee} & Fourth moment, selected density, 109/120 proposal\\[3pt]
968 & \texttt{14f544fb4293} & All fixed numerator moments; parallel 10/11 proposal\\[3pt]
969 & \texttt{654278118846} & Off-central fourth mean; 49/54 proposal\\[3pt]
970 & \texttt{026f70fc91f2} & Archived 68/75 proposal; not independently extracted or re-reviewed\\[3pt]
971 & \texttt{3f1d2a07631e} & Five ramified classes; local-domain improvement\\[3pt]
974 & \texttt{89c37b6707cb} & Parallel interpolation; 9081/10000 and 49/54; local domain\\[3pt]
977 & \texttt{1b473843c920} & Proposed $271/300$ and $19/21$ comparisons.\\[3pt]
989 & \texttt{0162599b1373} & Two mollifiers and capped reciprocal moments.\\[3pt]
991 & \texttt{b4d97d837214} & Fixed-shift extraction.\\[3pt]
993 & \texttt{5569defd8d77} & Prime-bias argument and the prime-power integer.\\[3pt]
994 & \texttt{c93ca65a510b} & Signed M\"obius sectors, order-$20$ Ising model, and theta arithmetic.\\[3pt]
996 & \texttt{56ec8747ccf7} & Fixed-field rational coefficient comparison.\\[3pt]
998 & \texttt{65a0e7e4016a} & An invertible signed average in one field.\\[3pt]
1000 & \texttt{6c4d5309d649} & Positive determinant and regularization.\\[3pt]
1001 & \texttt{f6beb3774ae7} & Quadratic-family limits and correlated strip models.\\[3pt]
1002 & \texttt{1319373815ee} & Uniform field average, filters, compactness, and larger masses.\\[3pt]
1003 & \texttt{669b6a01f4d4} & Finite order-$26$ constructions and quadratic-family estimates.\\[3pt]
1004 & \texttt{8e7375cd7590} & The failed $0.87$ comparison.\\[3pt]
1005 & \texttt{d3b3d19573b8} & Conditional all-height corollary; no new analytic estimate\\[3pt]
1006 & \texttt{0c838d130f27} & Related quartic filters and an explicitly unproved mixed mean\\[3pt]
1008 & \texttt{c1de6363f957} & Fixed-order ambiguity, not an actual off-critical zeta zero\\[3pt]
\end{longtable}
\normalsize

The new historical summaries for notes~957--974 use their detailed published statements and the cited source-identification notes. They are not fresh audits of the full analytic arguments. In particular, the archived $68/75$ work is included as a historical proposal, not as a newly reconstructed proof. The accounts of Liu and Xu--Sang use the exact versions recorded in notes~1002 and~1005 because fresh retrieval of their primary bodies was unsuccessful. These limitations do not affect the elementary identities proved directly in the paper; they do matter when judging the imported zero-free conclusions.

\subsection{Related results and their place in the argument}
The following related results are useful companions to the proofs above. They are summarized here, not used as substitutes for the open estimates in Section~\ref{sec:verification}.

\paragraph{Fixed-shift extraction.}
The source \cite{R991} studies the actual deformed coefficients
$a_h(\mathfrak I)=\sum_{\mathfrak A\mathfrak B=\mathfrak I}\mu_F(\mathfrak A)(N\mathfrak B)^{-h}$ and a complete centered Gaussian source $B(h,X)$. Its removable endpoint at $h=0$ is the fixed prime-power discrepancy. The stated Chebyshev interpolation uses $O_A(\log X)$ positive shifts and $O_A(\log^2X)$ total raw coefficient mass to recover that endpoint with error $O_A(X^{-A})$. A fixed real-shift average over an interval bounded away from zero retains every target zero on a positive-measure set of shifts. The source does not prove the subpower upper estimate for that average. This is an extraction result, not an imported inverse-$L$ moment bound.

\paragraph{Same-field causal replica capture.}
Over the Eisenstein field, \cite{R998} averages actual Euler-deleted native sources over all exterior ideals of norm at most $B$. The finite signed dilation kernel tends to $\chi(d)/N\rad(d)$ in a weighted absolute coefficient norm. Its limiting local series is
\[
 \frac{1-(1-1/q)z}{1-z},\qquad q=Np,
\]
with inverse coefficients $-q^{-1}(1-q^{-1})^{j-1}\chi(p)^j$ at $p^j$. The source proves finite-horizon norm equivalence with constants uniform in the displayed root-norm order. The complete limiting covariance has weight $1/N\rad(de)$, not the product of separate radical weights when primes overlap. Deletion is allowed for a small fraction of whole replica trajectories, not a different exceptional set at each physical time. No upper bound on the original native source is implied by invertibility of this averaging operation.

\paragraph{A complete signed divisibility contribution.}
The critical arithmetic section of \cite{R994} computes an infinite divisibility-sector coefficient with all tails retained and certifies $-9/5<L_{\rm div}<-17/10$. Its negative sign is a statement about that full specified sector, not the sign or an upper bound for the complete energy. The complementary nondivisibility contribution remains. The same source explicitly records that its primitive-ratio double limiting series is not absolutely summable, so an unrestricted interchange of the two limits is not available.


\paragraph{Quartic filters over the Gaussian field.}
The later note~1006 \cite{R1006} develops fourth-power filters over $\Q(\sqrt{-1})$ as a possible alternative to the sextic construction. Its reported retained squared norm is $1/\zeta_F^S(3)$, whereas a naive squareclass selector has squared norm
\[
 \prod_{\mathfrak p}\frac{N\mathfrak p+1}{2N\mathfrak p},
\]
which tends to zero. This identifies a loss in that particular selector, not an impossibility theorem for quartic methods. A separate comparison in the note would yield the boundary $437/500=0.874$ from a stated mixed mean of size $U^{6/5+\eps}$. That mean is unproved. It must not be listed as an achieved $0.874$ region. This paragraph reports the published scope; the full quartic proof and its computation collection were not re-audited for this revision.

\paragraph{Finite probability models with different zero geometries.}
Note~1008 \cite{R1008} also gives a probability-law counterpart to Proposition~\ref{prop:finite-insertion}. It is based on the finite order-$20$ Ising result, not on an order-$26$ probability realization. Its examples retain the stated positivity, strip, and leading growth properties while allowing nonreal zeros beyond a chosen height. The zero-counting error is $O(\sqrt T)$, not the sharper error for zeta. These examples delimit what the finite data establish; they are not examples of off-critical zeta zeros.

\subsection{Review and submission}
The manuscript and its source are prepared for review and compilation with standard LaTeX. No arXiv submission is made by publishing them in the repository. Before submission, authorship and contribution credit should be approved, and the analytic adaptations should receive independent mathematical review. The private motivating note discussed in \cite{R993} is not redistributed, and no authenticated public attribution or license is inferred for it.

\input{main.bbl}
\end{document}
